% ============================================================================
% MÓDULO 8. REPORTES AUTOMÁTICOS CON R MARKDOWN
% ============================================================================
\chapter{Reportes automáticos con R Markdown}
\label{cap:reportes}

\begin{objetivos}
\begin{itemize}
  \item Explicar por qué un \termino{reporte reproducible} ahorra horas y
        evita errores frente al reporte hecho a mano en Excel o PowerPoint.
  \item Crear un documento R Markdown (\archivo{.Rmd}) con encabezado YAML,
        texto en Markdown, bloques de código y código en línea, y convertirlo
        a HTML, PDF y Word con el botón \menu{Knit}.
  \item Controlar qué se ve en el reporte con las opciones de los bloques
        (\codigo{echo}, \codigo{include}, \codigo{fig.cap},
        \codigo{results = "asis"}, etc.).
  \item Presentar tablas profesionales con \codigo{knitr::kable()} y
        escribir \termino{texto dinámico} que cambia solo según los datos.
  \item Construir \termino{reportes parametrizados} (uno por tienda, por
        mes, por región) y generarlos todos con un script que registra lo
        que salió bien y lo que falló.
  \item Programar la ejecución automática (cron, Programador de tareas de
        Windows), enviar los reportes por correo y exportar reportes a Excel
        con varias hojas.
  \item Dar el salto a Quarto y aplicar buenas prácticas de reproducibilidad.
\end{itemize}
\end{objetivos}

Es el primer lunes del mes. Andrea, analista de la cadena TechRetail México,
llega a la oficina sabiendo lo que le espera: abrir el sistema, exportar las
ventas del mes a Excel, armar diez tablas dinámicas (una por sucursal), copiar
los totales a una presentación de PowerPoint, pegar las gráficas, escribir a
mano frases como ``las ventas crecieron 12\,\% respecto al mes anterior'' y,
por último, enviar diez correos a diez gerentes. Le toma casi dos días. Y el
mes pasado, al copiar y pegar, el gerente de la Sucursal Norte recibió las
cifras de la Sucursal Sur. Nadie lo notó hasta la junta de resultados.

Ahora imagina otra versión de ese lunes: Andrea llega, abre su correo y ya
tiene la confirmación de que a las 7:00 de la mañana un script de R generó
los diez reportes (con tablas, gráficas y conclusiones escritas en español a
partir de los datos) y se los envió a cada gerente. Andrea dedica su mañana a
lo que de verdad aporta valor: investigar por qué la Sucursal Centro cayó
tanto en diciembre.

Esa segunda versión es lo que vas a construir en este módulo. Usarás todo lo
que aprendiste de manipulación de datos (Módulo~\ref{cap:manipulacion}) y de
visualización (Módulo~\ref{cap:visualizacion}), y le agregarás una pieza
nueva: \paquete{R Markdown}, la herramienta que convierte tu análisis en un
documento terminado que se vuelve a generar con un clic, o sin ningún clic.

% ----------------------------------------------------------------------------
\section{El problema: el reporte mensual hecho a mano}
\label{sec:m8-problema}

Casi todas las empresas tienen un reporte que se hace ``cada mes, igual que
el mes pasado'': el reporte de ventas, el de cartera vencida, el de
inventario, el de rotación de personal. Cuando se hace a mano, el proceso
suele verse así:

\begin{enumerate}
  \item Exportar datos de un sistema a Excel.
  \item Limpiar, filtrar y resumir con fórmulas o tablas dinámicas.
  \item Copiar los resultados a Word o PowerPoint.
  \item Rehacer las gráficas o pegarlas como imagen.
  \item Escribir los comentarios (``crecimos'', ``caímos'', ``destaca'').
  \item Revisar, exportar a PDF y enviar.
\end{enumerate}

Cada paso manual es una oportunidad de error: una fila que no se seleccionó,
una fórmula que apunta a la celda equivocada, una gráfica del mes pasado que
nadie actualizó, un porcentaje escrito a mano que ya no corresponde a la
tabla. Y el costo en tiempo se multiplica: si el reporte es por tienda, son
diez veces el mismo trabajo.

\begin{table}[H]
\centering
\caption{Reporte manual frente a reporte reproducible}
\label{tab:m8-manual-vs-reproducible}
\small
\begin{tabularx}{\textwidth}{l X X}
\toprule
\textbf{Aspecto} & \textbf{Manual (Excel + PowerPoint)} &
  \textbf{Reproducible (R Markdown)} \\
\midrule
Tiempo por reporte & Horas, cada mes & Minutos la primera vez; segundos
  después \\
Errores de copiar y pegar & Frecuentes y difíciles de detectar & No hay
  copiar y pegar: los números salen del código \\
Diez versiones (una por tienda) & Diez veces el trabajo & Un bucle de
  diez vueltas \\
Texto de conclusiones & Se escribe a mano & Se genera con reglas
  (\codigo{if}/\codigo{else}) \\
Auditoría (``¿de dónde salió este número?'') & Difícil & El código lo
  documenta \\
Cambiar el formato (PDF, Word, HTML) & Rehacer & Cambiar una línea \\
\bottomrule
\end{tabularx}
\end{table}

\begin{concepto}{Reporte reproducible}
Un \termino{reporte reproducible} es un documento que se genera
\emph{completamente} a partir de los datos y de un código, sin pasos
manuales. Si los datos cambian (llega un mes nuevo, se corrige una venta),
basta con volver a generar el documento para que todos los números, tablas,
gráficas y frases se actualicen. Cualquier persona con los mismos datos y el
mismo código obtiene exactamente el mismo reporte.
\end{concepto}

\subsection{Qué es R Markdown y cómo funciona}

\termino{R Markdown} es un formato de archivo (extensión \archivo{.Rmd}) que
mezcla tres cosas en un solo documento:

\begin{itemize}
  \item \textbf{Texto} escrito en \termino{Markdown}, un lenguaje muy
        sencillo para dar formato (títulos, negritas, listas) con símbolos
        como \codigo{\#} o \codigo{**}.
  \item \textbf{Código R} en bloques llamados \termino{chunks}, cuyos
        resultados (tablas, gráficas, números) se insertan en el documento.
  \item Un \textbf{encabezado} en formato \termino{YAML} con el título, el
        autor y el formato de salida.
\end{itemize}

Cuando presionas el botón \menu{Knit} (``tejer'') en RStudio, ocurren dos
pasos, que resume la figura~\ref{fig:m8-flujo}. Primero, el paquete
\paquete{knitr} ejecuta cada bloque de código y ``teje'' sus resultados con
el texto, generando un archivo Markdown puro (\archivo{.md}). Después,
\termino{pandoc}, un convertidor universal de documentos que viene incluido
con RStudio, transforma ese \archivo{.md} en el formato final: una página
HTML, un documento de Word o, pasando por LaTeX, un PDF. El paquete
\paquete{rmarkdown} coordina todo el proceso con su función
\codigo{render()}.

\begin{figure}[H]
\centering
\begin{tikzpicture}[
  archivo/.style={draw=primario, fill=primario!6, rounded corners=3pt,
                  minimum width=1.6cm, minimum height=1cm, align=center,
                  font=\sffamily\footnotesize},
  motor/.style={draw=acento, fill=acento!10, ellipse, minimum width=1.5cm,
                minimum height=0.9cm, align=center,
                font=\sffamily\footnotesize\bfseries},
  salida/.style={draw=exito, fill=exito!8, rounded corners=3pt,
                 minimum width=1.7cm, minimum height=0.75cm, align=center,
                 font=\sffamily\footnotesize},
  flecha/.style={-{Stealth[length=2.5mm]}, thick, color=gristexto}]
  \node[archivo] (rmd) {\textbf{.Rmd}\\texto + código};
  \node[motor, right=1.1cm of rmd] (knitr) {knitr};
  \node[archivo, right=0.6cm of knitr] (md) {\textbf{.md}\\texto +\\resultados};
  \node[motor, right=1.3cm of md] (pandoc) {pandoc};
  \node[salida, right=0.8cm of pandoc, yshift=1.05cm] (html) {HTML};
  \node[salida, right=0.8cm of pandoc] (docx) {Word (.docx)};
  \node[salida, right=0.8cm of pandoc, yshift=-1.05cm] (tex) {LaTeX};
  \node[salida, right=0.4cm of tex] (pdf) {PDF};
  \draw[flecha] (rmd) -- node[above, font=\sffamily\scriptsize]{ejecuta} (knitr);
  \draw[flecha] (knitr) -- (md);
  \draw[flecha] (md) -- node[above, font=\sffamily\scriptsize]{convierte} (pandoc);
  \draw[flecha] (pandoc.east) -- (html.west);
  \draw[flecha] (pandoc.east) -- (docx.west);
  \draw[flecha] (pandoc.east) -- (tex.west);
  \draw[flecha] (tex) -- (pdf);
  \node[font=\sffamily\scriptsize, color=gristexto, below=0.35cm of md]
    {\codigo{rmarkdown::render()} coordina todo el proceso (botón Knit)};
\end{tikzpicture}
\caption{Cómo se convierte un archivo R Markdown en un reporte terminado.}
\label{fig:m8-flujo}
\end{figure}

Lo importante de este flujo es que \textbf{el reporte nunca se edita a mano}:
si algo está mal, corriges el \archivo{.Rmd} (o los datos) y vuelves a
tejer. El documento final es un producto desechable; lo valioso es el
\archivo{.Rmd}.

Para que veas cómo trabaja \paquete{knitr} sin salir de la consola, puedes
pedirle que teja un pequeño texto con \codigo{knitr::knit(text = ...)}. Cada
elemento del vector es una línea del documento:

\begin{rcode}
library(tidyverse)   # dplyr, readr, ggplot2, tidyr, purrr...
library(knitr)       # knit() y kable()
library(scales)      # dollar(), percent(), comma()

# Un mini documento R Markdown guardado en un vector de texto
rmd_mini <- c(
  "Este año vendimos `r format(1234567, big.mark = ',')` pesos.",
  "",
  "```{r suma}",
  "sum(1:10)",
  "```"
)
# knitr ejecuta el código y devuelve Markdown puro (lo que ve pandoc)
cat(knit(text = rmd_mini, quiet = TRUE))
\end{rcode}
\begin{rsalida}
Este año vendimos 1,234,567 pesos.


```r
sum(1:10)
```

```
## [1] 55
```
\end{rsalida}

Observa lo que pasó: la expresión \codigo{`r format(...)`} de la primera
línea desapareció y en su lugar quedó el número ya formateado
(\codigo{1,234,567}). El bloque de código se conservó (para que se vea en el
reporte) y debajo apareció su resultado, precedido de \codigo{\#\#}. Eso es
exactamente lo que knitr hace con un reporte completo; después pandoc le da
formato bonito.

\begin{consejo}
Para seguir este módulo necesitas los paquetes \paquete{rmarkdown} y
\paquete{knitr}. RStudio los instala automáticamente la primera vez que
creas un archivo R Markdown, pero también puedes instalarlos tú. Para los
PDF necesitarás además LaTeX (sección~\ref{sec:m8-formatos}).
\end{consejo}

\begin{rnoejecutar}
install.packages(c("rmarkdown", "knitr", "writexl", "readxl"))
\end{rnoejecutar}

% ----------------------------------------------------------------------------
\section{Anatomía de un documento R Markdown}
\label{sec:m8-anatomia}

\subsection{Crear tu primer .Rmd en RStudio}

\begin{enumerate}
  \item Ve a \menu{File > New File > R Markdown\ldots}
  \item Escribe un título (por ejemplo, ``Mi primer reporte de ventas'') y tu
        nombre, deja seleccionado \menu{HTML} y presiona \menu{OK}.
  \item RStudio abre un documento de ejemplo. Guárdalo con
        \tecla{Ctrl} + \tecla{S} dentro de la carpeta del curso (la que
        contiene \archivo{datasets/}), por ejemplo como
        \archivo{reportes/01-primer-reporte.Rmd}.
  \item Presiona el botón \menu{Knit} (el ovillo de estambre azul) o el
        atajo \tecla{Ctrl} + \tecla{Shift} + \tecla{K}. En unos segundos se
        abre la vista previa del HTML y el archivo
        \archivo{01-primer-reporte.html} aparece junto al \archivo{.Rmd}.
\end{enumerate}

En la carpeta \archivo{codigo/reportes/} del libro encontrarás los archivos
completos de este módulo, listos para copiar a tu carpeta \archivo{reportes/}:
\archivo{01-primer-reporte.Rmd}, \archivo{02-reporte-ejecutivo.Rmd},
\archivo{generar\_reportes.R} y \archivo{reporte\_excel.R}. La estructura de
carpetas que usaremos es:

\begin{textocodigo}
curso-bi-rstudio/              <- directorio de trabajo (proyecto RStudio)
|-- datasets/                  <- datos de entrada (CSV)
|   |-- transacciones.csv
|   |-- tiendas.csv
|   `-- ...
|-- reportes/                  <- plantillas .Rmd y scripts
|   |-- 01-primer-reporte.Rmd
|   |-- 02-reporte-ejecutivo.Rmd
|   |-- generar_reportes.R
|   `-- reporte_excel.R
`-- resultados/                <- salidas (se pueden borrar y regenerar)
    |-- reporte_ventas.xlsx
    `-- reportes/
        |-- 2023-12_T01_sucursal-centro_2024-01-05.html
        `-- ...
\end{textocodigo}

Un archivo \archivo{.Rmd} tiene tres tipos de contenido, que muestra la
figura~\ref{fig:m8-anatomia}. Vamos a recorrer el archivo
\archivo{01-primer-reporte.Rmd} por partes.

\begin{figure}[H]
\centering
\begin{tikzpicture}[
  zona/.style={draw, rounded corners=2pt, minimum width=7.2cm, align=left,
               font=\ttfamily\scriptsize, inner sep=5pt, text width=6.9cm},
  nota/.style={font=\sffamily\footnotesize, align=left, text width=5.6cm},
  flecha/.style={-{Stealth[length=2mm]}, thick, color=gristexto}]
  \node[zona, draw=morado, fill=morado!6] (yaml)
    {-{}-{}-\\title: "Mi primer reporte de ventas"\\output: html\_document\\-{}-{}-};
  \node[zona, draw=gristexto, fill=gray!5, below=0.25cm of yaml] (texto)
    {\#\# Introducción\\Este es mi **primer reporte automático**.};
  \node[zona, draw=secundario, fill=secundario!6, below=0.25cm of texto] (chunk)
    {```\{r tabla-trimestres\}\\ventas \%>\% group\_by(trimestre) \%>\% ...\\```};
  \node[zona, draw=acento, fill=acento!8, below=0.25cm of chunk] (inline)
    {El archivo tiene `r nrow(ventas)` ventas.};
  \node[nota, right=0.5cm of yaml] (n1)
    {\textbf{\color{morado}1. Encabezado YAML}: título, autor, fecha y
     formato de salida, entre dos líneas \texttt{-{}-{}-}.};
  \node[nota, right=0.5cm of texto] (n2)
    {\textbf{2. Texto en Markdown}: lo que se lee en el reporte.};
  \node[nota, right=0.5cm of chunk] (n3)
    {\textbf{\color{secundario}3. Bloque de código (chunk)}: R se ejecuta y
     su resultado se inserta.};
  \node[nota, right=0.5cm of inline] (n4)
    {\textbf{\color{acento}4. Código en línea}: un valor calculado dentro
     de una frase.};
  \draw[flecha] (n1.west) -- (yaml.east);
  \draw[flecha] (n2.west) -- (texto.east);
  \draw[flecha] (n3.west) -- (chunk.east);
  \draw[flecha] (n4.west) -- (inline.east);
\end{tikzpicture}
\caption{Las cuatro piezas de un documento R Markdown.}
\label{fig:m8-anatomia}
\end{figure}

\subsection{El encabezado YAML}

El documento empieza con un bloque entre dos líneas de tres guiones. Es el
\termino{encabezado YAML} (\emph{YAML Ain't Markup Language}), un formato de
configuración del tipo \codigo{clave: valor}:

\begin{textocodigo}
---
title: "Mi primer reporte de ventas"
author: "Tu nombre"
date: "`r format(Sys.Date(), '%d/%m/%Y')`"
output: html_document
---
\end{textocodigo}

\begin{itemize}
  \item \codigo{title}, \codigo{author} y \codigo{date} aparecen en la
        portada del reporte. Observa que la fecha es código R: cada vez que
        tejes, se escribe la fecha de ese día.
  \item \codigo{output} indica el formato: \codigo{html\_document},
        \codigo{pdf\_document} o \codigo{word\_document}.
  \item Los textos con caracteres especiales (dos puntos, comillas, acentos
        al inicio) van entre comillas.
\end{itemize}

\begin{advertencia}
En YAML \textbf{la sangría importa}: las opciones de un formato se escriben
debajo de él con dos espacios de sangría, y nunca con tabuladores. Un
espacio de más o de menos rompe el documento
(sección~\ref{sec:m8-errores}).
\end{advertencia}

\subsection{Texto en Markdown}

Después del encabezado escribes el texto del reporte. Markdown usa símbolos
sencillos para dar formato; la tabla~\ref{tab:m8-markdown} resume los que
usarás el 95\,\% del tiempo.

\begin{table}[H]
\centering
\caption{Sintaxis básica de Markdown}
\label{tab:m8-markdown}
\small
\begin{tabularx}{\textwidth}{l l X}
\toprule
\textbf{Quieres\ldots} & \textbf{Escribes} & \textbf{Notas} \\
\midrule
Título de nivel 1, 2, 3 & \codigo{\# Título}, \codigo{\#\# Sección},
  \codigo{\#\#\# Subsección} & Deja un espacio después de \codigo{\#} y una
  línea en blanco antes \\
Negritas & \codigo{**texto**} & También \codigo{\_\_texto\_\_} \\
Cursivas & \codigo{*texto*} & También \codigo{\_texto\_} \\
Código en línea & \codigo{`read\_csv()`} & Letra de máquina de escribir \\
Lista con viñetas & \codigo{- elemento} & Una línea en blanco antes de la
  lista \\
Lista numerada & \codigo{1. elemento} & La numeración se ajusta sola \\
Enlace & \codigo{[texto](https://...)} & \\
Imagen & \codigo{![pie](logo.png)} & Ruta relativa al \archivo{.Rmd} \\
Línea horizontal & \codigo{-{}-{}-} & En una línea sola \\
Tabla & \codigo{| A | B |} y \codigo{|---|---|} & Mejor generarlas con
  \codigo{kable()} \\
Salto de párrafo & Una línea en blanco & Un solo salto de línea no separa
  párrafos \\
\bottomrule
\end{tabularx}
\end{table}

Un ejemplo de texto con varios de estos elementos:

\begin{textocodigo}
## Introducción

Este es mi **primer reporte automático**. Analiza las ventas de 2023 de la
cadena *TechRetail México*. Puntos clave:

- Los datos vienen de `ventas_retail.csv`.
- Se actualiza cada vez que presiono **Knit**.

| Trimestre | Meta       |
|-----------|-----------:|
| Q1        | $1,000,000 |
\end{textocodigo}

\subsection{Bloques de código (chunks)}

Un \termino{chunk} empieza con tres acentos graves y \codigo{\{r\}}, y
termina con tres acentos graves. Opcionalmente lleva un nombre
(\codigo{cargar-datos}) y opciones separadas por comas. En RStudio lo
insertas con \tecla{Ctrl} + \tecla{Alt} + \tecla{I} y lo ejecutas, sin
tejer, con la flechita verde de su esquina.

\begin{textocodigo}
```{r cargar-datos}
library(dplyr)
library(readr)
library(ggplot2)

# El .Rmd se ejecuta en SU carpeta. Buscamos datasets/ ahí mismo
# o en la carpeta superior (si el .Rmd está en curso/reportes/).
carpeta <- if (dir.exists("datasets")) "datasets" else "../datasets"
ventas <- read_csv(file.path(carpeta, "ventas_retail.csv"),
                   show_col_types = FALSE)
```
\end{textocodigo}

Fíjate en la búsqueda de la carpeta de datos. Cuando tejes, R usa como
directorio de trabajo \textbf{la carpeta donde está guardado el
\archivo{.Rmd}}, no la de tu proyecto. Si el \archivo{.Rmd} vive en
\archivo{reportes/}, la ruta \archivo{datasets/ventas\_retail.csv} no existe
desde ahí: hay que subir un nivel (\archivo{../datasets}). Volveremos a esto
en la sección~\ref{sec:m8-directorio}, porque es el error número uno de quien
empieza con R Markdown.

\subsection{Código en línea}

Para meter un número calculado \emph{dentro de una frase} se usa el
\termino{código en línea}: un acento grave, la letra \codigo{r}, un espacio,
la expresión y otro acento grave.

\begin{textocodigo}
El archivo tiene `r nrow(ventas)` ventas y el total vendido en el año
fue de **$`r format(round(sum(ventas$total)), big.mark = ",")`**.
\end{textocodigo}

Antes de ponerlo en el reporte, conviene probar la expresión en la consola.
Carguemos los datos igual que lo hace el reporte (con el directorio de
trabajo en la carpeta del curso):

\begin{rcode}
ventas <- read_csv("datasets/ventas_retail.csv", show_col_types = FALSE)

nrow(ventas)                                   # número de ventas
format(round(sum(ventas$total)), big.mark = ",")  # total con comas
\end{rcode}
\begin{rsalida}
[1] 365
[1] "513,917"
\end{rsalida}

En el reporte, la frase se leerá: ``El archivo tiene 365 ventas y el total
vendido en el año fue de \textbf{\$513,917}''. Si mañana llegan más
datos, la frase se corrige sola.

\subsection{El documento completo y el botón Knit}

Juntando las piezas, este es el resto de \archivo{01-primer-reporte.Rmd}
(después del encabezado, el texto y el bloque de carga):

\begin{textocodigo}
```{r setup, include=FALSE}
# Opciones para todos los bloques: mostrar código, ocultar mensajes
knitr::opts_chunk$set(echo = TRUE, message = FALSE, warning = FALSE,
                      fig.width = 7, fig.height = 3.5)
```

## Ventas por trimestre

```{r tabla-trimestres}
ventas %>%
  group_by(trimestre) %>%
  summarise(ventas = round(sum(total)), n_ventas = n()) %>%
  knitr::kable(col.names = c("Trimestre", "Ventas ($)", "Núm. de ventas"),
               format.args = list(big.mark = ","))
```

## Tendencia mensual

```{r grafica-mensual, fig.cap = "Ventas totales por mes, 2023"}
ventas %>%
  group_by(mes) %>%
  summarise(ventas = sum(total)) %>%
  ggplot(aes(x = mes, y = ventas, group = 1)) +
  geom_line(color = "#2a78d6", linewidth = 0.8) +
  geom_point(color = "#2a78d6") +
  scale_y_continuous(labels = scales::dollar) +
  labs(x = NULL, y = "Ventas") +
  theme_minimal() +
  theme(panel.grid.minor = element_blank(),
        axis.text.x = element_text(angle = 45, hjust = 1))
```
\end{textocodigo}

Al presionar \menu{Knit}, en el panel \menu{Render} de RStudio verás cómo
knitr procesa cada bloque (\codigo{1/9}, \codigo{2/9}\ldots) y al final la
línea \codigo{Output created: 01-primer-reporte.html}. El HTML resultante es
un solo archivo autocontenido (las gráficas van incrustadas): lo puedes
enviar por correo y se abre en cualquier navegador, sin necesidad de R.

\begin{hazlotu}
Crea \archivo{reportes/01-primer-reporte.Rmd} con el contenido anterior y
tejelo. Después agrega una sección \codigo{\#\# Ventas por día de la semana}
con una tabla que use \codigo{dia\_semana}, y vuelve a tejer. ¿Cuánto tardó
en actualizarse todo el documento?
\end{hazlotu}

% ----------------------------------------------------------------------------
\section{Opciones de los bloques de código}
\label{sec:m8-opciones}

Un reporte para el director no debe mostrar código, mensajes de carga de
paquetes ni advertencias; un tutorial para el equipo de análisis sí debe
mostrar el código. Eso se controla con las \termino{opciones de chunk}, que
se escriben dentro de las llaves después del nombre:

\begin{textocodigo}
```{r grafica-ventas, echo = FALSE, fig.cap = "Ventas por mes",
    fig.width = 7, fig.height = 3.5, out.width = "90%"}
\end{textocodigo}

La tabla~\ref{tab:m8-opciones} reúne las opciones que más usarás.

\begin{table}[H]
\centering
\caption{Opciones de chunk más usadas en reportes de negocio}
\label{tab:m8-opciones}
\small
\begin{tabularx}{\textwidth}{l l X}
\toprule
\textbf{Opción} & \textbf{Valor típico} & \textbf{Qué hace} \\
\midrule
\codigo{echo} & \codigo{FALSE} & Oculta el código; muestra el resultado \\
\codigo{eval} & \codigo{FALSE} & Muestra el código pero no lo ejecuta \\
\codigo{include} & \codigo{FALSE} & Ejecuta el código pero no muestra
  nada (ni código ni resultado). Ideal para cargar datos \\
\codigo{message} & \codigo{FALSE} & Oculta mensajes (p.~ej., al cargar
  paquetes) \\
\codigo{warning} & \codigo{FALSE} & Oculta advertencias \\
\codigo{error} & \codigo{TRUE} & Muestra el error y sigue tejiendo (por
  defecto, un error detiene todo) \\
\codigo{fig.cap} & \codigo{"Título"} & Pie de la figura (acepta una
  variable de R) \\
\codigo{fig.width}, \codigo{fig.height} & \codigo{7}, \codigo{3.5} &
  Tamaño de la gráfica en pulgadas \\
\codigo{out.width} & \codigo{"90\%"} & Ancho con el que se muestra en el
  documento \\
\codigo{fig.align} & \codigo{"center"} & Alineación de la figura \\
\codigo{results} & \codigo{"asis"} & El texto que imprime el código se
  trata como Markdown, no como salida de consola \\
\codigo{results} & \codigo{"hide"} & Oculta la salida impresa \\
\codigo{cache} & \codigo{TRUE} & Guarda el resultado y no vuelve a
  ejecutar el bloque si su código no cambió \\
\bottomrule
\end{tabularx}
\end{table}

Veamos la diferencia entre \codigo{echo = FALSE} e \codigo{include = FALSE}
tejiendo un pequeño documento en la consola:

\begin{rcode}
doc_opciones <- c(
  "```{r a, echo = FALSE}",
  "sum(1:10)",
  "```",
  "",
  "```{r b, include = FALSE}",
  "total <- sum(1:10)",
  "```",
  "",
  "El total es `r total`."
)
cat(knit(text = doc_opciones, quiet = TRUE))
\end{rcode}
\begin{rsalida}

```
## [1] 55
```



El total es 55.
\end{rsalida}

Con \codigo{echo = FALSE} el código desapareció pero quedó su resultado
(\codigo{\#\# [1] 55}). Con \codigo{include = FALSE} no quedó nada visible,
pero el objeto \codigo{total} sí se creó, y por eso la frase final puede
usarlo. Ese es el patrón de un reporte ejecutivo: bloques invisibles que
calculan, y texto que muestra.

\subsection{El chunk de configuración global}

Repetir \codigo{echo = FALSE, message = FALSE} en cada bloque es tedioso.
Lo habitual es un primer bloque, llamado \codigo{setup} o
\codigo{configuracion}, con \codigo{include = FALSE}, que fija las opciones
\emph{por defecto} para todo el documento con
\codigo{knitr::opts\_chunk\$set()}:

\begin{textocodigo}
```{r configuracion, include=FALSE}
knitr::opts_chunk$set(
  echo = FALSE,      # un reporte ejecutivo NO muestra código
  message = FALSE,   # sin mensajes de carga de paquetes
  warning = FALSE,   # sin advertencias
  fig.width = 7, fig.height = 3.4, out.width = "90%",
  fig.align = "center"
)
library(dplyr)
library(ggplot2)
```
\end{textocodigo}

Cualquier bloque puede cambiar una opción para sí mismo: si en un reporte
con \codigo{echo = FALSE} quieres mostrar un bloque, le pones
\codigo{echo = TRUE} solo a ese.

\subsection{\texorpdfstring{\codigo{results = "asis"}}{results = "asis"}:
cuando el código escribe el reporte}

Por defecto, lo que tu código imprime con \codigo{cat()} o \codigo{print()}
aparece en un recuadro gris, como salida de consola. Con
\codigo{results = "asis"} (``tal cual''), ese texto se inserta como parte del
documento, así que puedes \emph{generar Markdown desde R}: listas, títulos o
párrafos que dependen de los datos.

\begin{rcode}
doc_asis <- c(
  "```{r lista, echo = FALSE, results = 'asis'}",
  "tiendas_top <- c('Sucursal Outlet', 'Sucursal Sur')",
  "cat(paste0('- **', tiendas_top, '**'), sep = '\\n')",
  "```"
)
cat(knit(text = doc_asis, quiet = TRUE))
\end{rcode}
\begin{rsalida}
- **Sucursal Outlet**
- **Sucursal Sur**
\end{rsalida}

El resultado ya no lleva \codigo{\#\#} ni recuadro: es una lista de Markdown
que pandoc mostrará con viñetas y negritas. En el caso práctico usaremos esta
técnica para escribir las conclusiones automáticas del reporte.

\subsection{\texorpdfstring{\codigo{cache}}{cache}: no recalcular lo que no
cambió}

Si un bloque tarda mucho (leer millones de filas, entrenar un modelo del
Módulo~\ref{cap:ml}), \codigo{cache = TRUE} guarda su resultado en una
carpeta \archivo{*\_cache/} y en el siguiente Knit lo reutiliza si el código
del bloque no cambió. Úsalo con cuidado: si cambian los \emph{datos} pero no
el código, knitr no se entera y te mostrará resultados viejos. En reportes
mensuales automáticos, donde lo que cambia son precisamente los datos, lo
más seguro es no usar caché en los bloques que leen datos.

\begin{hazlotu}
En tu \archivo{01-primer-reporte.Rmd} cambia el bloque \codigo{setup} para
que \codigo{echo = FALSE} y vuelve a tejer. Luego agrega
\codigo{fig.cap} con un título que incluya el año y cambia
\codigo{out.width} a \codigo{"70\%"}. Compara ambas versiones.
\end{hazlotu}

% ----------------------------------------------------------------------------
\section{\texorpdfstring{Tablas profesionales con \codigo{knitr::kable()}}{Tablas profesionales con knitr::kable()}}
\label{sec:m8-kable}

Si imprimes un tibble en un reporte, se ve como salida de consola: letra de
máquina, \codigo{<dbl>} y \codigo{<chr>} en los encabezados, números sin
separador de miles. La función \codigo{kable()} del paquete \paquete{knitr}
convierte un data frame en una tabla con formato de documento (HTML, LaTeX o
Word, según la salida).

\begin{ejemplo}{Ventas de diciembre por tienda}
El director quiere una tabla con las ventas de diciembre de 2023 de cada
tienda: nombre, ciudad, número de transacciones, ventas y ticket promedio,
ordenada de mayor a menor. Preparamos los datos con \paquete{dplyr} y
después le damos formato con \codigo{kable()}.
\end{ejemplo}

\begin{rcode}
# Datos que usaremos en el resto del módulo
transacciones <- read_csv("datasets/transacciones.csv",
                          show_col_types = FALSE)
tiendas   <- read_csv("datasets/tiendas.csv", show_col_types = FALSE)
productos <- read_csv("datasets/productos.csv", show_col_types = FALSE)

# Resumen de diciembre por tienda
ventas_dic <- transacciones %>%
  filter(mes == "2023-12") %>%
  group_by(tienda_id) %>%
  summarise(transacciones = n(),
            ventas = sum(total_transaccion),
            ticket = mean(total_transaccion), .groups = "drop") %>%
  left_join(tiendas %>% select(tienda_id, nombre_tienda, ciudad),
            by = "tienda_id") %>%
  select(nombre_tienda, ciudad, transacciones, ventas, ticket) %>%
  arrange(desc(ventas))

ventas_dic
\end{rcode}
\begin{rsalida}
# A tibble: 10 x 5
   nombre_tienda    ciudad      transacciones ventas ticket
   <chr>            <chr>               <int>  <dbl>  <dbl>
 1 Sucursal Outlet  Tijuana                12 16554.  1379.
 2 Sucursal Sur     Monterrey               8 13317.  1665.
 3 Sucursal Express Querétaro              10 11194.  1119.
 4 Sucursal Plaza A CDMX                    9 10945.  1216.
 5 Sucursal Norte   Guadalajara            15 10049.   670.
 6 Sucursal Este    CDMX                   11  9983.   908.
 7 Sucursal Plaza B Guadalajara            11  9928.   903.
 8 Sucursal Centro  CDMX                    5  5020.  1004.
 9 Sucursal Premium CDMX                    5  4916.   983.
10 Sucursal Oeste   Puebla                  5  4488.   898.
\end{rsalida}

Así se ve en la consola: útil para ti, pero no para un director. Ahora con
\codigo{kable()}. Usamos \codigo{format = "pipe"} para que veas en el libro
la tabla Markdown que genera; dentro de un \archivo{.Rmd} no hace falta
ponerlo, porque knitr elige el formato adecuado para HTML, PDF o Word.

\begin{rcode}
kable(ventas_dic,
      format = "pipe",                        # Markdown (para verla aquí)
      col.names = c("Tienda", "Ciudad", "Transacciones",
                    "Ventas ($)", "Ticket ($)"),  # encabezados legibles
      digits = 0,                             # sin decimales
      format.args = list(big.mark = ","),     # separador de miles
      align = c("l", "l", "r", "r", "r"),     # texto izq., números der.
      caption = "Ventas de diciembre de 2023 por tienda")
\end{rcode}
\begin{rsalida}


Table: Ventas de diciembre de 2023 por tienda

|Tienda           |Ciudad      | Transacciones| Ventas ($)| Ticket ($)|
|:----------------|:-----------|-------------:|----------:|----------:|
|Sucursal Outlet  |Tijuana     |            12|     16,554|      1,379|
|Sucursal Sur     |Monterrey   |             8|     13,317|      1,665|
|Sucursal Express |Querétaro   |            10|     11,194|      1,119|
|Sucursal Plaza A |CDMX        |             9|     10,945|      1,216|
|Sucursal Norte   |Guadalajara |            15|     10,049|        670|
|Sucursal Este    |CDMX        |            11|      9,983|        908|
|Sucursal Plaza B |Guadalajara |            11|      9,928|        903|
|Sucursal Centro  |CDMX        |             5|      5,020|      1,004|
|Sucursal Premium |CDMX        |             5|      4,916|        983|
|Sucursal Oeste   |Puebla      |             5|      4,488|        898|
\end{rsalida}

Cada argumento resuelve un problema de presentación:

\begin{itemize}
  \item \codigo{col.names} cambia los nombres técnicos
        (\codigo{nombre\_tienda}) por nombres para humanos.
  \item \codigo{digits} redondea todas las columnas numéricas y
        \codigo{format.args = list(big.mark = ",")} agrega separador de
        miles.
  \item \codigo{align} alinea: \codigo{"l"} izquierda para texto,
        \codigo{"r"} derecha para números (así se comparan las cifras por
        su magnitud), \codigo{"c"} centro.
  \item \codigo{caption} pone el título numerado (``Tabla 1: \ldots'').
\end{itemize}

Los dos puntos en la segunda línea de la tabla Markdown
(\codigo{|-----:|}) son los que indican a pandoc la alineación a la derecha.

\subsection{Formato de negocio antes de la tabla}

\codigo{kable()} no sabe poner el signo de pesos ni porcentajes. Cuando
quieras ese formato, conviértelo tú antes con las funciones de
\paquete{scales} (las mismas que usaste en los ejes de ggplot2 en el
Módulo~\ref{cap:visualizacion}). Ojo: al formatear, la columna se vuelve
texto, así que ordena \emph{antes} de formatear.

\begin{rcode}
tabla_dic <- ventas_dic %>%
  mutate(participacion = ventas / sum(ventas),       # % del total
         ventas        = dollar(ventas, accuracy = 1),
         ticket        = dollar(ticket, accuracy = 0.01),
         participacion = percent(participacion, accuracy = 0.1)) %>%
  select(-ciudad) %>%
  slice_head(n = 5)                                  # top 5

kable(tabla_dic, format = "pipe", align = c("l", "r", "r", "r", "r"),
      col.names = c("Tienda", "Transacciones", "Ventas", "Ticket",
                    "% del total"))
\end{rcode}
\begin{rsalida}


|Tienda           | Transacciones|  Ventas|    Ticket| % del total|
|:----------------|-------------:|-------:|---------:|-----------:|
|Sucursal Outlet  |            12| $16,554| $1,379.48|       17.2%|
|Sucursal Sur     |             8| $13,317| $1,664.67|       13.8%|
|Sucursal Express |            10| $11,194| $1,119.44|       11.6%|
|Sucursal Plaza A |             9| $10,945| $1,216.07|       11.4%|
|Sucursal Norte   |            15| $10,049|   $669.96|       10.4%|
\end{rsalida}

La tabla dice que en diciembre la Sucursal Outlet fue la que más vendió
(\$16,554, el 17.2\,\% de la empresa) y que las cinco primeras tiendas tienen
tickets promedio muy distintos: la Sucursal Sur logró \$13,317 con solo
8 transacciones (ticket de \$1,664.67), mientras la Sucursal Norte necesitó
15 transacciones para vender \$10,049. Esa es la clase de observación que
luego puedes convertir en texto dinámico.

\subsection{Más allá de kable: kableExtra y gt}

Para tablas con colores, filas resaltadas, encabezados agrupados o notas al
pie, existen dos paquetes muy usados (no los necesitas para este módulo):

\begin{rnoejecutar}
# kableExtra: añade estilo a kable() (HTML y PDF)
install.packages("kableExtra")
library(kableExtra)
kable(tabla_dic) %>%
  kable_styling(bootstrap_options = c("striped", "hover")) %>%
  row_spec(1, bold = TRUE, background = "#fde6dc")   # resalta la 1a fila

# gt: "gramática de tablas" de Posit, muy flexible
install.packages("gt")
library(gt)
ventas_dic %>%
  gt() %>%
  fmt_currency(columns = c(ventas, ticket)) %>%
  tab_header(title = "Ventas de diciembre de 2023")
\end{rnoejecutar}

\begin{hazlotu}
Haz una tabla con \codigo{kable()} de las ventas totales de 2023 por
\codigo{metodo\_pago} (de \archivo{transacciones.csv}), con las ventas en
formato de pesos, el porcentaje del total y el título ``Ventas por método de
pago''.
\end{hazlotu}

% ----------------------------------------------------------------------------
\section{Texto dinámico: frases que se escriben solas}
\label{sec:m8-texto}

Un reporte ejecutivo no es solo tablas: es \emph{narrativa}. ``Las ventas
crecieron 19.7\,\% respecto a noviembre''. Si esa frase se escribe a mano,
el mes siguiente puede quedarse con el número viejo, o peor, decir
``crecieron'' cuando en realidad cayeron. El \termino{texto dinámico} es
texto que R construye a partir de los datos: el número \emph{y} la palabra
(``crecieron''/``cayeron'') salen del cálculo.

\subsection{Formatear números para leerlos en una frase}

\begin{rcode}
x <- 96393.537
dollar(x)                       # pesos, con separador de miles
dollar(x, accuracy = 1)         # sin centavos
dollar(x, scale = 1e-3, suffix = " mil", accuracy = 0.1)
comma(1234567)                  # separador de miles
percent(0.19731, accuracy = 0.1)
format(x, big.mark = ",", nsmall = 2)          # con base R
format(as.Date("2023-12-15"), "%d de %B de %Y")  # fecha en texto
\end{rcode}
\begin{rsalida}
[1] "$96,393.54"
[1] "$96,394"
[1] "$96.4 mil"
[1] "1,234,567"
[1] "19.7%"
[1] "96,393.54"
[1] "15 de diciembre de 2023"
\end{rsalida}

La última línea usa el idioma de tu computadora para el nombre del mes. Si
tu Windows está en inglés, diría \codigo{December}. Para que el reporte
salga siempre en español, en el caso práctico usaremos un vector propio con
los nombres de los meses.

\subsection{Frases que cambian según los datos}

\begin{ejemplo}{La frase de crecimiento mensual}
Queremos una función que reciba las ventas del mes actual y del anterior y
devuelva una frase correcta en cualquier caso: si crecieron, si cayeron, si
se quedaron igual o si no hay dato del mes anterior.
\end{ejemplo}

\begin{rcode}
frase_variacion <- function(actual, anterior, periodo = "el mes anterior") {
  # Caso 1: no hay con qué comparar
  if (is.na(anterior) || anterior == 0) {
    return(paste0("Las ventas fueron de ", dollar(actual, accuracy = 1),
                  "; no hay datos para comparar con ", periodo, "."))
  }
  var <- actual / anterior - 1                   # variación relativa
  # Caso 2, 3 y 4: la palabra depende del signo
  verbo <- if (abs(var) < 0.005) {
    "se mantuvieron sin cambio"
  } else if (var > 0) {
    paste("crecieron", percent(var, accuracy = 0.1))
  } else {
    paste("cayeron", percent(abs(var), accuracy = 0.1))
  }
  paste0("Las ventas ", verbo, " en comparación con ", periodo, " (de ",
         dollar(anterior, accuracy = 1), " a ",
         dollar(actual, accuracy = 1), ").")
}

# Probamos los cuatro casos
frase_variacion(12000, 10000)
frase_variacion(8000, 10000)
frase_variacion(10020, 10000)
frase_variacion(5000, NA)
\end{rcode}
\begin{rsalida}
[1] "Las ventas crecieron 20.0% en comparación con el mes anterior (de $10,000 a $12,000)."
[1] "Las ventas cayeron 20.0% en comparación con el mes anterior (de $10,000 a $8,000)."
[1] "Las ventas se mantuvieron sin cambio en comparación con el mes anterior (de $10,000 a $10,020)."
[1] "Las ventas fueron de $5,000; no hay datos para comparar con el mes anterior."
\end{rsalida}

Fíjate en el tercer caso: una variación de 0.2\,\% se reporta como ``sin
cambio'', porque decir ``crecieron 0.2\,\%'' exagera algo que es ruido. Estas
reglas de redacción son decisiones de negocio que conviene acordar con quien
lee el reporte.

Ahora apliquemos la función a los datos reales: ventas totales de la
empresa en diciembre contra noviembre.

\begin{rcode}
ventas_mes <- transacciones %>%
  group_by(mes) %>%
  summarise(ventas = sum(total_transaccion), .groups = "drop")

actual   <- ventas_mes$ventas[ventas_mes$mes == "2023-12"]
anterior <- ventas_mes$ventas[ventas_mes$mes == "2023-11"]
frase_variacion(actual, anterior, periodo = "noviembre")
\end{rcode}
\begin{rsalida}
[1] "Las ventas crecieron 19.7% en comparación con noviembre (de $80,507 a $96,394)."
\end{rsalida}

En el \archivo{.Rmd}, esta frase se inserta con código en línea:

\begin{textocodigo}
## Resumen

`r frase_variacion(actual, anterior, periodo = "noviembre")`
\end{textocodigo}

\subsection{Una frase para cada tienda}

Como la función trabaja con dos números, podemos aplicarla a las diez
tiendas con \codigo{map2\_chr()} de \paquete{purrr} (que recorre dos vectores
a la vez y devuelve texto):

\begin{rcode}
comparacion <- transacciones %>%
  filter(mes %in% c("2023-11", "2023-12")) %>%
  group_by(tienda_id, mes) %>%
  summarise(ventas = sum(total_transaccion), .groups = "drop") %>%
  pivot_wider(names_from = mes, values_from = ventas) %>%
  rename(nov = `2023-11`, dic = `2023-12`) %>%
  left_join(tiendas %>% select(tienda_id, nombre_tienda),
            by = "tienda_id") %>%
  mutate(frase = map2_chr(dic, nov, frase_variacion,
                          periodo = "noviembre"))

# Imprimimos las frases de las 3 primeras tiendas
cat(paste0(comparacion$nombre_tienda[1:3], ": ",
           comparacion$frase[1:3]), sep = "\n")
\end{rcode}
\begin{rsalida}
Sucursal Centro: Las ventas cayeron 52.6% en comparación con noviembre (de $10,593 a $5,020).
Sucursal Norte: Las ventas crecieron 25.8% en comparación con noviembre (de $7,990 a $10,049).
Sucursal Sur: Las ventas crecieron 97.4% en comparación con noviembre (de $6,746 a $13,317).
\end{rsalida}

Tres tiendas, tres historias distintas, cero frases escritas a mano. La
Sucursal Centro cayó a menos de la mitad y la Sucursal Sur casi duplicó sus
ventas. Pero cuidado con la interpretación: cada tienda tiene solo entre 3 y
16 transacciones al mes en estos datos, así que variaciones tan grandes
pueden deberse al azar. Un buen reporte automático también advierte eso
(lo haremos en las conclusiones del caso práctico).

\subsection{Conclusiones automáticas con reglas}

Un paso más allá de la frase es la \termino{conclusión automática}: una
regla de negocio que decide qué mensaje dar. Por ejemplo: si la variación es
mayor a 10\,\%, ``crecimiento sólido''; si cae más de 10\,\%, ``alerta''. Con
\codigo{case\_when()} lo aplicamos a todas las tiendas:

\begin{rcode}
alertas <- comparacion %>%
  mutate(variacion = dic / nov - 1,
         estado = case_when(
           variacion >  0.10 ~ "Crecimiento sólido",
           variacion >= 0    ~ "Estable",
           variacion > -0.10 ~ "Ligera caída",
           TRUE              ~ "Alerta"
         )) %>%
  arrange(variacion) %>%
  select(nombre_tienda, variacion, estado)

alertas %>%
  mutate(variacion = percent(variacion, accuracy = 0.1)) %>%
  kable(format = "pipe", col.names = c("Tienda", "Variación", "Estado"))
\end{rcode}
\begin{rsalida}


|Tienda           |Variación |Estado             |
|:----------------|:---------|:------------------|
|Sucursal Centro  |-52.6%    |Alerta             |
|Sucursal Oeste   |-41.6%    |Alerta             |
|Sucursal Premium |-37.3%    |Alerta             |
|Sucursal Outlet  |-13.5%    |Alerta             |
|Sucursal Norte   |25.8%     |Crecimiento sólido |
|Sucursal Express |34.2%     |Crecimiento sólido |
|Sucursal Plaza B |84.2%     |Crecimiento sólido |
|Sucursal Sur     |97.4%     |Crecimiento sólido |
|Sucursal Plaza A |170.1%    |Crecimiento sólido |
|Sucursal Este    |264.4%    |Crecimiento sólido |
\end{rsalida}

Y con \codigo{results = "asis"}, una lista de Markdown para el director que
solo menciona las tiendas en alerta:

\begin{rcode}
en_alerta <- alertas %>% filter(estado == "Alerta")

if (nrow(en_alerta) == 0) {
  cat("- Ninguna tienda cayó más de 10 % este mes.\n")
} else {
  cat(paste0("- **", en_alerta$nombre_tienda, "**: cayó ",
             percent(abs(en_alerta$variacion), accuracy = 0.1),
             " contra noviembre."), sep = "\n")
}
\end{rcode}
\begin{rsalida}
- **Sucursal Centro**: cayó 52.6% contra noviembre.
- **Sucursal Oeste**: cayó 41.6% contra noviembre.
- **Sucursal Premium**: cayó 37.3% contra noviembre.
- **Sucursal Outlet**: cayó 13.5% contra noviembre.
\end{rsalida}

\begin{negocio}{El reporte que avisa solo}
En muchas empresas el reporte mensual es un documento de 30 páginas que
nadie lee completo. Un reporte con conclusiones automáticas pone arriba, en
tres líneas, lo que requiere atención (``3 tiendas en alerta'') y deja el
detalle para quien lo necesite. El analista deja de ``escribir el reporte''
y pasa a \emph{diseñar las reglas} con las que se escribe, que es un trabajo
mucho más valioso.
\end{negocio}

\begin{hazlotu}
Modifica \codigo{frase\_variacion()} para que, cuando el crecimiento sea
mayor a 50\,\%, agregue al final ``Verifica que no haya ventas
extraordinarias.''. Pruébala con \codigo{frase\_variacion(16000, 10000)}.
\end{hazlotu}

% ----------------------------------------------------------------------------
\section{Formatos de salida: HTML, PDF y Word}
\label{sec:m8-formatos}

El mismo \archivo{.Rmd} puede producir tres formatos distintos. Solo cambia
el campo \codigo{output} del encabezado.

\subsection{HTML}

\begin{textocodigo}
output:
  html_document:
    toc: true            # índice (table of contents)
    toc_float: true      # el índice flota a la izquierda
    theme: flatly        # tema visual (cerulean, cosmo, journal...)
    code_folding: hide   # el código se oculta con un botón "Code"
    number_sections: true
\end{textocodigo}

El HTML es el formato más flexible: admite gráficas interactivas (como las
de \paquete{plotly} del Módulo~\ref{cap:shiny}), índice flotante y el botón
\codigo{code\_folding}, que es ideal cuando el lector \emph{a veces} quiere
ver el código: está oculto, pero a un clic.

\subsection{PDF}

\begin{textocodigo}
output:
  pdf_document:
    toc: true
    fig_caption: true
    extra_dependencies: ["float"]   # paquete LaTeX adicional
geometry: margin=2cm                # márgenes de la página
lang: es-MX                         # idioma (guiones, "Figura")
\end{textocodigo}

Para producir PDF, R Markdown necesita una instalación de
\termino{LaTeX}, el sistema de composición tipográfica con el que se hacen
libros y artículos (este libro, por ejemplo). Si no tienes LaTeX, la forma
más sencilla es \paquete{TinyTeX}, una distribución pequeña que se instala
desde R:

\begin{rnoejecutar}
install.packages("tinytex")
tinytex::install_tinytex()   # descarga e instala LaTeX (una sola vez)
tinytex::is_tinytex()        # TRUE si quedó instalado
\end{rnoejecutar}

La primera vez que tejas un PDF, TinyTeX descargará automáticamente los
paquetes de LaTeX que falten; por eso esa primera vez tarda más y necesita
Internet.

\subsection{Word}

\begin{textocodigo}
output:
  word_document:
    toc: true
    fig_caption: true
    reference_docx: plantilla_corporativa.docx
\end{textocodigo}

Word es el formato preferido cuando alguien más va a \emph{editar} el
reporte (agregar comentarios, ajustar un párrafo). La opción
\codigo{reference\_docx} es muy poderosa: le das un documento de Word cuyo
estilo copiará tu reporte (fuentes, colores de títulos, márgenes, encabezado
con el logo de la empresa). Para crear esa plantilla:

\begin{enumerate}
  \item Teje una vez cualquier \archivo{.Rmd} a Word.
  \item Abre el \archivo{.docx} generado y modifica los \emph{estilos}
        (no el texto): clic derecho en ``Título 1'' $\rightarrow$
        \menu{Modificar}, y así con ``Normal'', ``Título 2'', etc.
  \item Guárdalo junto al \archivo{.Rmd} con un nombre como
        \archivo{plantilla\_corporativa.docx} y agrega la opción
        \codigo{reference\_docx}.
\end{enumerate}

\subsection{Varios formatos en un mismo archivo}

Puedes declarar los tres formatos en el encabezado (es lo que hace nuestro
reporte ejecutivo). El botón \menu{Knit} muestra entonces un menú para
elegir, y desde R puedes generar uno o todos:

\begin{rnoejecutar}
library(rmarkdown)
render("reportes/02-reporte-ejecutivo.Rmd",
       output_format = "pdf_document")   # solo PDF
render("reportes/02-reporte-ejecutivo.Rmd",
       output_format = "all")            # HTML, PDF y Word
\end{rnoejecutar}

\begin{table}[H]
\centering
\caption{¿Qué formato de salida elegir?}
\label{tab:m8-formatos}
\small
\begin{tabularx}{\textwidth}{l X X X}
\toprule
 & \textbf{HTML} & \textbf{PDF} & \textbf{Word} \\
\midrule
Requiere & Nada extra & LaTeX (TinyTeX) & Nada extra (Word para
  abrirlo) \\
Ideal para & Reportes en pantalla, intranet, correo & Documentos formales,
  impresión, archivo & Documentos que otros van a editar \\
Interactividad & Sí (plotly, tablas DT, pestañas) & No & No \\
Aspecto & Controlado con \codigo{theme} y CSS & Muy profesional, controlado
  con LaTeX & Controlado con \codigo{reference\_docx} \\
Riesgo & Algunos correos bloquean HTML adjunto & Errores de LaTeX al
  instalar & Se puede modificar el contenido \\
\bottomrule
\end{tabularx}
\end{table}

\begin{consejo}
Una opción que dependa del formato se puede decidir desde el propio código.
\codigo{knitr::pandoc\_to()} devuelve el formato que se está generando
(\codigo{"html"}, \codigo{"latex"}, \codigo{"docx"}). Nuestro reporte
ejecutivo la usa para no pedir \codigo{fig.align} en Word, que no la admite:
\codigo{es\_word <- knitr::pandoc\_to("docx")}.
\end{consejo}

% ----------------------------------------------------------------------------
\section{Reportes parametrizados}
\label{sec:m8-parametros}

Hasta ahora nuestro reporte habla de ``la empresa''. Pero el gerente de la
Sucursal Norte quiere \emph{su} reporte, y la directora quiere el de
diciembre, no el de noviembre. Copiar el \archivo{.Rmd} diez veces y cambiar
el número de tienda en cada copia sería volver al trabajo manual. La
solución son los \termino{parámetros}.

\begin{concepto}{Reporte parametrizado}
Un \termino{reporte parametrizado} es una plantilla \archivo{.Rmd} que
declara en su encabezado unos \emph{parámetros} (por ejemplo,
\codigo{tienda} y \codigo{mes}) con un valor por defecto. Dentro del código
se usan como \codigo{params\$tienda} y \codigo{params\$mes}. Al generar el
reporte puedes darles otro valor, y obtienes un documento distinto desde la
misma plantilla.
\end{concepto}

\subsection{Declarar los parámetros en el YAML}

\begin{textocodigo}
params:
  tienda:
    label: "Tienda (1 a 10)"
    value: 1
    input: select
    choices: [1, 2, 3, 4, 5, 6, 7, 8, 9, 10]
  mes:
    label: "Mes del reporte (AAAA-MM)"
    value: "2023-12"
    input: text
  ruta_datos:
    label: "Carpeta de datos"
    value: "datasets"
    input: text
\end{textocodigo}

La forma mínima es \codigo{tienda: 1}. La forma extendida (con
\codigo{label}, \codigo{value}, \codigo{input}) sirve para el botón
\menu{Knit > Knit with Parameters\ldots}: RStudio abre un pequeño formulario
con una lista desplegable para la tienda y una caja de texto para el mes.
Así, alguien que no programa puede generar el reporte de su tienda eligiendo
en un menú.

\subsection{Usar los parámetros en el código}

Dentro del \archivo{.Rmd} existe automáticamente una lista llamada
\codigo{params}. Se usa como cualquier lista de R:

\begin{textocodigo}
```{r datos}
id_tienda  <- as.integer(params$tienda)
mes_actual <- params$mes
ventas_mes <- ventas %>%
  filter(tienda_id == id_tienda, mes == mes_actual)
```

# Reporte de la tienda `r id_tienda` (`r mes_actual`)
\end{textocodigo}

Un truco muy útil para \emph{desarrollar} la plantilla: en la consola crea
tú una lista \codigo{params} con los valores que quieras probar, y ejecuta
los bloques uno por uno. El código se comporta igual que al tejer. Hagamos
eso con la validación que usa nuestro reporte ejecutivo:

\begin{rcode}
# En la consola simulamos los parámetros que recibiría el .Rmd
params <- list(tienda = 3, mes = "2023-12")

id_tienda  <- as.integer(params$tienda)
mes_actual <- params$mes

# Validaciones: mejor un error claro que un reporte vacío
info_tienda <- tiendas %>% filter(tienda_id == id_tienda)
if (nrow(info_tienda) == 0) stop("La tienda ", id_tienda, " no existe.")
if (!mes_actual %in% transacciones$mes) {
  stop("No hay datos del mes ", mes_actual)
}

# Mes anterior: el día 1 del mes menos un día cae en el mes previo
mes_previo <- format(as.Date(paste0(mes_actual, "-01")) - 1, "%Y-%m")

info_tienda$nombre_tienda
mes_previo
\end{rcode}
\begin{rsalida}
[1] "Sucursal Sur"
[1] "2023-11"
\end{rsalida}

El cálculo del mes anterior merece una explicación: \codigo{"2023-12-01"}
menos un día es \codigo{"2023-11-30"}, y su formato \codigo{"\%Y-\%m"} es
\codigo{"2023-11"}. El truco funciona también en enero (el 1 de enero menos
un día es diciembre del año anterior), algo en lo que suelen fallar las
fórmulas hechas a mano.

\begin{advertencia}
Dentro de un documento que se está tejiendo, \codigo{params} es de solo
lectura: \codigo{params\$tienda <- 5} produce el error \emph{``no se puede
cambiar el valor de un vínculo bloqueado para `params'\,''}. Si necesitas
transformar un parámetro, guárdalo en otra variable, como hicimos con
\codigo{id\_tienda}.
\end{advertencia}

\subsection{Generar el reporte con otros parámetros desde R}

La función \codigo{rmarkdown::render()} recibe los parámetros en una lista.
Todo lo que no indiques toma el valor por defecto del YAML:

\begin{rnoejecutar}
library(rmarkdown)
# Reporte de la Sucursal Sur (tienda 3), diciembre de 2023, en PDF
render("reportes/02-reporte-ejecutivo.Rmd",
       output_format = "pdf_document",
       output_file   = "reporte_T03_2023-12.pdf",
       output_dir    = "resultados/reportes",
       params        = list(tienda = 3, mes = "2023-12"),
       knit_root_dir = getwd())
\end{rnoejecutar}

Los argumentos de \codigo{render()} que usarás siempre son:

\begin{itemize}
  \item \codigo{input}: la plantilla \archivo{.Rmd}.
  \item \codigo{output\_format}: \codigo{"html\_document"},
        \codigo{"pdf\_document"}, \codigo{"word\_document"} o \codigo{"all"}.
  \item \codigo{output\_file} y \codigo{output\_dir}: nombre y carpeta del
        archivo generado (si no los das, se guarda junto al \archivo{.Rmd}
        con el mismo nombre, y cada reporte \emph{sobrescribe} al anterior).
  \item \codigo{params}: la lista de parámetros.
  \item \codigo{knit\_root\_dir}: el directorio de trabajo en el que se
        ejecutará el código del \archivo{.Rmd} (lo explicamos enseguida).
  \item \codigo{envir}: el entorno donde se crean los objetos;
        \codigo{new.env()} garantiza que cada reporte empiece limpio.
  \item \codigo{quiet = TRUE}: no imprime el avance de knitr.
\end{itemize}

\subsection{¿Desde qué carpeta se ejecuta el .Rmd?}
\label{sec:m8-directorio}

Este detalle causa más dolores de cabeza que cualquier otro, así que vale la
pena entenderlo bien. Cuando tejes, \textbf{el directorio de trabajo es la
carpeta donde está el \archivo{.Rmd}}, sin importar cuál sea el directorio
de trabajo de tu consola o de tu proyecto. Con nuestra estructura:

\begin{itemize}
  \item Tu consola está en \archivo{curso-bi-rstudio/}: ahí
        \archivo{datasets/transacciones.csv} existe.
  \item El \archivo{.Rmd} está en \archivo{curso-bi-rstudio/reportes/}: desde
        ahí, \archivo{datasets/transacciones.csv} \emph{no} existe (habría
        que escribir \archivo{../datasets/transacciones.csv}).
\end{itemize}

Hay tres soluciones, y nuestro reporte ejecutivo combina las dos primeras
para funcionar en cualquier caso:

\begin{enumerate}
  \item \textbf{\codigo{knit\_root\_dir} en \codigo{render()}}: le dice a
        knitr que ejecute el código desde otra carpeta. Con
        \codigo{knit\_root\_dir = getwd()} (o la ruta del proyecto) el código
        corre en la carpeta del curso y \archivo{datasets/} se encuentra.
        En RStudio puedes lograr lo mismo desde el botón Knit:
        \menu{Knit > Knit Directory > Project Directory}.
  \item \textbf{Un parámetro con la ruta y una búsqueda robusta}: el
        parámetro \codigo{ruta\_datos} indica la carpeta de datos y el código
        la busca en el directorio actual y, si no, en el superior:
\end{enumerate}

\begin{textocodigo}
```{r datos}
candidatas <- c(params$ruta_datos, file.path("..", params$ruta_datos))
carpeta_datos <- candidatas[dir.exists(candidatas)][1]
if (is.na(carpeta_datos)) {
  stop("No encontré la carpeta de datos '", params$ruta_datos,
       "'. Directorio actual: ", getwd(),
       ". Usa knit_root_dir en render() o el parámetro ruta_datos.")
}
transacciones <- read_csv(file.path(carpeta_datos, "transacciones.csv"),
                          show_col_types = FALSE)
```
\end{textocodigo}

\begin{enumerate}[resume]
  \item \textbf{Paquete \paquete{here}}: \codigo{here::here("datasets",
        "transacciones.csv")} construye la ruta desde la raíz del proyecto
        de RStudio, esté donde esté el \archivo{.Rmd}. Es una excelente
        opción si trabajas con proyectos (\archivo{.Rproj}).
\end{enumerate}

Así, el reporte funciona si presionas Knit en RStudio (se ejecuta en
\archivo{reportes/} y encuentra \archivo{../datasets}), si lo generas con
\codigo{render(..., knit\_root\_dir = getwd())} (encuentra
\archivo{datasets}), y si alguien tiene los datos en otra carpeta (le pasa
\codigo{params = list(ruta\_datos = "D:/datos")}). Y si no los encuentra,
el mensaje de error dice exactamente qué pasó y cómo arreglarlo.

\begin{hazlotu}
Genera desde la consola el reporte ejecutivo de la Sucursal Premium
(tienda 9) para octubre de 2023 en Word. Luego ábrelo con
\menu{Knit > Knit with Parameters\ldots} y genera el de la tienda 4 de junio
en HTML.
\end{hazlotu}

% ----------------------------------------------------------------------------
\section{Automatización: un reporte por tienda, sin tocar nada}
\label{sec:m8-automatizacion}

Con una plantilla parametrizada, generar los diez reportes es cuestión de un
bucle. El script \archivo{generar\_reportes.R} hace cuatro cosas:
construye los nombres de archivo, llama a \codigo{render()} una vez por
tienda, atrapa los errores para que una tienda con problemas no detenga a
las demás, y deja un registro de lo generado. Veamos cada pieza por separado
y luego el script completo.

\subsection{Nombres de archivo con fecha}

Un buen nombre de archivo dice qué es, de qué periodo y cuándo se generó, y
se ordena bien alfabéticamente. Usaremos el formato
\archivo{AAAA-MM\_Tnn\_nombre-de-tienda\_AAAA-MM-DD.html}. Como los nombres
de las tiendas tienen espacios (y podrían tener acentos), los
``limpiamos'':

\begin{rcode}
# "Sucursal Plaza A" -> "sucursal-plaza-a" (sin acentos ni espacios)
limpiar_nombre <- function(x) {
  x <- chartr("áéíóúüñÁÉÍÓÚÜÑ", "aeiouunAEIOUUN", x)  # quita acentos
  x <- gsub("[^A-Za-z0-9]+", "-", x)     # lo que no sea letra -> guion
  tolower(gsub("^-|-$", "", x))          # sin guiones al inicio/final
}
limpiar_nombre(c("Sucursal Plaza A", "Tienda Querétaro Ñuñoa"))

# En el script real la fecha es Sys.Date(); aquí la fijamos para que
# el ejemplo dé siempre el mismo resultado
fecha_ejecucion <- as.Date("2024-01-05")
mes_reporte <- "2023-12"

archivos <- sprintf("%s_T%02d_%s_%s.html", mes_reporte,
                    tiendas$tienda_id,
                    limpiar_nombre(tiendas$nombre_tienda),
                    format(fecha_ejecucion, "%Y-%m-%d"))
head(archivos, 4)
\end{rcode}
\begin{rsalida}
[1] "sucursal-plaza-a"       "tienda-queretaro-nunoa"
[1] "2023-12_T01_sucursal-centro_2024-01-05.html"
[2] "2023-12_T02_sucursal-norte_2024-01-05.html" 
[3] "2023-12_T03_sucursal-sur_2024-01-05.html"   
[4] "2023-12_T04_sucursal-este_2024-01-05.html"  
\end{rsalida}

\codigo{sprintf()} arma el texto con un molde: \codigo{\%s} inserta un texto
y \codigo{\%02d} un número entero con dos dígitos (\codigo{1} $\rightarrow$
\codigo{01}), para que la tienda 10 quede después de la 9 al ordenar.

\subsection{La tabla de trabajos: qué reportes generar}

Cuando los reportes dependen de dos parámetros (tienda y mes, o tienda y
formato), conviene armar primero una \termino{tabla de trabajos} con todas
las combinaciones y después recorrerla. \codigo{expand\_grid()} de
\paquete{tidyr} crea todas las combinaciones:

\begin{rcode}
trabajos <- expand_grid(tienda_id = tiendas$tienda_id,
                        mes = c("2023-11", "2023-12")) %>%
  left_join(tiendas %>% select(tienda_id, nombre_tienda),
            by = "tienda_id") %>%
  mutate(archivo = sprintf("%s_T%02d_%s.html", mes, tienda_id,
                           limpiar_nombre(nombre_tienda)))

nrow(trabajos)
head(trabajos, 3)

# La lista de parámetros de la primera fila, lista para render()
parametros <- list(tienda = trabajos$tienda_id[1],
                   mes = trabajos$mes[1])
str(parametros)
\end{rcode}
\begin{rsalida}
[1] 20
# A tibble: 3 x 4
  tienda_id mes     nombre_tienda   archivo                         
      <dbl> <chr>   <chr>           <chr>                           
1         1 2023-11 Sucursal Centro 2023-11_T01_sucursal-centro.html
2         1 2023-12 Sucursal Centro 2023-12_T01_sucursal-centro.html
3         2 2023-11 Sucursal Norte  2023-11_T02_sucursal-norte.html 
List of 2
 $ tienda: num 1
 $ mes   : chr "2023-11"
\end{rsalida}

\subsection{Errores controlados con tryCatch}

Si al generar el reporte de la tienda 7 ocurre un error (por ejemplo, no
tiene ventas ese mes), sin protección el bucle se detiene y las tiendas 8, 9
y 10 se quedan sin reporte. \codigo{tryCatch()} ejecuta un código y, si
falla, en vez de detenerse ejecuta la función que le indiques en
\codigo{error}. Para practicar sin generar archivos, simulemos
\codigo{render()} con una función que valida los parámetros igual que el
\archivo{.Rmd}:

\begin{rcode}
# Simula render(): falla igual que el .Rmd cuando algo no cuadra
generar_uno <- function(tienda, mes) {
  if (!tienda %in% tiendas$tienda_id) {
    stop("La tienda ", tienda, " no existe.")
  }
  if (!mes %in% transacciones$mes) stop("No hay datos del mes ", mes)
  sprintf("reporte_T%02d_%s.html", tienda, mes)
}

pendientes <- tibble(tienda = c(1, 2, 15, 3),
                     mes = c("2023-12", "2023-12", "2023-12", "2024-01"))
registro <- vector("list", nrow(pendientes))

for (fila in seq_len(nrow(pendientes))) {
  resultado <- tryCatch({
    archivo <- generar_uno(pendientes$tienda[fila], pendientes$mes[fila])
    paste("OK:", archivo)
  }, error = function(e) {
    paste("ERROR:", conditionMessage(e))   # el mensaje del error
  })
  registro[[fila]] <- tibble(tienda = pendientes$tienda[fila],
                             mes = pendientes$mes[fila],
                             estado = if (startsWith(resultado, "OK"))
                               "OK" else "ERROR",
                             detalle = resultado)
}
registro <- bind_rows(registro)
registro
\end{rcode}
\begin{rsalida}
# A tibble: 4 x 4
  tienda mes     estado detalle                            
   <dbl> <chr>   <chr>  <chr>                              
1      1 2023-12 OK     OK: reporte_T01_2023-12.html       
2      2 2023-12 OK     OK: reporte_T02_2023-12.html       
3     15 2023-12 ERROR  ERROR: La tienda 15 no existe.     
4      3 2024-01 ERROR  ERROR: No hay datos del mes 2024-01
\end{rsalida}

El bucle terminó las cuatro vueltas. Dos reportes salieron bien y dos
fallaron, pero en lugar de detenerse, cada error quedó anotado con su causa:
la tienda 15 no existe y enero de 2024 no tiene datos. Esa tabla es el
\termino{registro} (\emph{log}): al día siguiente, en vez de preguntarte
``¿se habrán generado los reportes?'', lo abres y lo sabes.

\subsection{El script completo}

Este es \archivo{generar\_reportes.R}. Como llama a \codigo{render()} y
escribe archivos, lo mostramos como referencia; en la sección de caso
práctico verás su ejecución real.

\begin{rnoejecutar}[title=reportes/generar\_reportes.R (parte 1: configuración)]
suppressPackageStartupMessages({
  library(rmarkdown)
  library(readr)
  library(dplyr)
})

# ---- 1. Configuración -----------------------------------------------
# Argumentos opcionales de la terminal: mes (AAAA-MM) y formato
argumentos  <- commandArgs(trailingOnly = TRUE)
mes_reporte <- if (length(argumentos) >= 1) argumentos[1] else "2023-12"
formato     <- if (length(argumentos) >= 2) argumentos[2] else "html"
# En producción usarías el mes anterior a hoy, por ejemplo:
# mes_reporte <- format(lubridate::floor_date(Sys.Date(), "month") - 1,
#                       "%Y-%m")

# Formatos disponibles: nombre corto -> formato de R Markdown y extensión
formatos <- list(
  html = c("html_document", "html"),
  pdf  = c("pdf_document", "pdf"),
  word = c("word_document", "docx")
)
if (!formato %in% names(formatos)) {
  stop("Formato no válido: ", formato, ". Usa html, pdf o word.")
}

# Carpeta del curso: la actual o la superior (si ejecutas desde reportes/)
raiz <- if (dir.exists("datasets")) "." else
        if (dir.exists("../datasets")) ".."
if (is.null(raiz)) {
  stop("No encuentro datasets/. Ejecuta desde la carpeta del curso.")
}
raiz <- normalizePath(raiz)

plantilla      <- file.path(raiz, "reportes", "02-reporte-ejecutivo.Rmd")
carpeta_salida <- file.path(raiz, "resultados", "reportes")
dir.create(carpeta_salida, recursive = TRUE, showWarnings = FALSE)
archivo_log    <- file.path(carpeta_salida, "registro_reportes.csv")
fecha_hoy      <- format(Sys.Date(), "%Y-%m-%d")
\end{rnoejecutar}

La primera parte lee dos argumentos opcionales con
\codigo{commandArgs(trailingOnly = TRUE)}: así el mismo script sirve para
cualquier mes y formato cuando se ejecuta desde la terminal
(\codigo{Rscript reportes/generar\_reportes.R 2023-12 pdf}). También
localiza la carpeta del curso y crea la carpeta de salida.

\begin{rnoejecutar}[title=reportes/generar\_reportes.R (parte 2: el bucle)]
# ---- 2. Lista de tiendas --------------------------------------------
tiendas <- read_csv(file.path(raiz, "datasets", "tiendas.csv"),
                    show_col_types = FALSE)
# (aquí va la función limpiar_nombre() que ya viste)

# ---- 3. Bucle: un reporte por tienda --------------------------------
cat("Generando", nrow(tiendas), "reportes del mes", mes_reporte,
    "en formato", formato, "\n")
registro <- vector("list", nrow(tiendas))

for (i in seq_len(nrow(tiendas))) {
  id     <- tiendas$tienda_id[i]
  nombre <- tiendas$nombre_tienda[i]
  archivo <- sprintf("%s_T%02d_%s_%s.%s", mes_reporte, id,
                     limpiar_nombre(nombre), fecha_hoy,
                     formatos[[formato]][2])
  inicio <- Sys.time()

  # tryCatch: si UNA tienda falla, se registra el error y seguimos
  resultado <- tryCatch({
    render(
      input         = plantilla,
      output_format = formatos[[formato]][1],
      output_file   = archivo,
      output_dir    = carpeta_salida,
      params        = list(tienda = id, mes = mes_reporte),
      knit_root_dir = raiz,       # el código del .Rmd corre en la raíz
      envir         = new.env(),  # entorno limpio para cada reporte
      quiet         = TRUE
    )
    "OK"
  }, error = function(e) {
    paste("ERROR:", conditionMessage(e))
  })

  segundos <- round(as.numeric(difftime(Sys.time(), inicio,
                                        units = "secs")), 1)
  registro[[i]] <- tibble(
    fecha_hora = format(Sys.time(), "%Y-%m-%d %H:%M:%S"),
    tienda_id  = id,
    tienda     = nombre,
    mes        = mes_reporte,
    archivo    = archivo,
    estado     = if (resultado == "OK") "OK" else "ERROR",
    detalle    = resultado,
    segundos   = segundos
  )
  cat(sprintf("  [%2d/%d] %-18s %-5s %s\n", i, nrow(tiendas), nombre,
              registro[[i]]$estado, archivo))
}

# ---- 4. Registro (log) ----------------------------------------------
registro <- bind_rows(registro)
# append = TRUE agrega al historial; los encabezados solo la primera vez
write_csv(registro, archivo_log, append = file.exists(archivo_log))

n_ok <- sum(registro$estado == "OK")
cat("\nListo:", n_ok, "reportes generados,", nrow(registro) - n_ok,
    "con error. Carpeta:", carpeta_salida, "\n")
cat("Registro guardado en:", archivo_log, "\n")
\end{rnoejecutar}

Tres detalles importantes del bucle:

\begin{itemize}
  \item \codigo{knit\_root\_dir = raiz} hace que el \archivo{.Rmd} encuentre
        \archivo{datasets/} (sección~\ref{sec:m8-directorio}).
  \item \codigo{envir = new.env()} da a cada reporte un entorno vacío. Sin
        esto, los objetos del reporte de la tienda 1 seguirían existiendo al
        generar el de la tienda 2, y un error de programación (usar una
        variable que no se recalculó) podría pasar desapercibido y
        \emph{mezclar datos de dos tiendas}: justo el error de copiar y
        pegar que queríamos eliminar.
  \item El registro se guarda en CSV con \codigo{append = TRUE}: cada
        ejecución agrega sus filas, así que con el tiempo tienes el historial
        completo de qué se generó y cuándo.
\end{itemize}

\subsection{Programar la ejecución}

Un script que genera todo con un comando ya es un gran avance. El último
paso es que se ejecute \emph{solo}, por ejemplo el día 1 de cada mes a las
7:00. Para eso se usa el programador de tareas del sistema operativo, que
llama a \termino{Rscript}, el programa que ejecuta un script de R desde la
terminal sin abrir RStudio.

\subsubsection{Linux y Mac: cron}

\termino{cron} es el programador de tareas de Linux y macOS. Se edita con
el comando \codigo{crontab -e} en la terminal, y cada línea es una tarea con
cinco campos de tiempo (minuto, hora, día del mes, mes, día de la semana)
seguidos del comando:

\begin{textocodigo}
# min hora dia mes dia_semana  comando
# Día 1 de cada mes a las 7:00: reportes del mes anterior en PDF
# (todo el comando va en UNA sola línea del crontab)
0 7 1 * * cd /home/andrea/curso-bi-rstudio && /usr/bin/Rscript reportes/generar_reportes.R $(date -d "last month" +\%Y-\%m) pdf >> resultados/cron.log 2>&1
\end{textocodigo}

\begin{itemize}
  \item \codigo{0 7 1 * *}: minuto 0, hora 7, día 1, cualquier mes,
        cualquier día de la semana. (\codigo{0 8 * * 1} sería ``cada lunes a
        las 8:00''.)
  \item \codigo{cd \ldots \&\&}: primero se cambia a la carpeta del curso,
        para que las rutas relativas funcionen.
  \item \codigo{\$(date -d "last month" +\%Y-\%m)} calcula el mes anterior
        (en cron el símbolo \codigo{\%} se escribe \codigo{\textbackslash\%}).
        En Mac, usa \codigo{date -v-1m +\%Y-\%m}. Otra opción es calcular el
        mes dentro del script (la línea comentada en la parte 1).
  \item \codigo{>> resultados/cron.log 2>\&1} guarda todo lo que imprime el
        script, incluidos los errores, en un archivo.
\end{itemize}

\subsubsection{Windows: Programador de tareas}

En Windows se usa el \termino{Programador de tareas} (\emph{Task
Scheduler}):

\begin{enumerate}
  \item Abre el menú Inicio y busca ``Programador de tareas''.
  \item Clic en \menu{Crear tarea básica\ldots}, ponle nombre (``Reportes
        mensuales'') y elige \menu{Mensualmente}, día 1, 7:00.
  \item En \menu{Acción} elige \menu{Iniciar un programa}.
  \item En \menu{Programa o script} escribe la ruta completa de
        \archivo{Rscript.exe}, que está en la carpeta \archivo{bin} de
        tu instalación de R (por ejemplo, dentro de
        \archivo{C:\textbackslash Program Files\textbackslash
        R\textbackslash R-4.3.3}).
  \item En \menu{Agregar argumentos} escribe
        \archivo{reportes\textbackslash generar\_reportes.R} y en
        \menu{Iniciar en} la carpeta del curso (por ejemplo
        \archivo{C:\textbackslash Users\textbackslash
        andrea\textbackslash curso-bi-rstudio}).
  \item Finaliza. Para probarla, clic derecho sobre la tarea $\rightarrow$
        \menu{Ejecutar}.
\end{enumerate}

Si prefieres hacerlo desde R, el paquete \paquete{taskscheduleR} crea la
tarea por ti (y RStudio le agrega un complemento con formulario en
\menu{Addins}). En Linux/Mac, el equivalente es \paquete{cronR}:

\begin{rnoejecutar}
install.packages("taskscheduleR")
library(taskscheduleR)
taskscheduler_create(
  taskname  = "reportes_mensuales",
  rscript   = file.path("C:/Users/andrea/curso-bi-rstudio",
                        "reportes/generar_reportes.R"),
  schedule  = "MONTHLY",
  days      = 1,
  starttime = "07:00"
)
taskscheduler_ls()                        # lista las tareas programadas
taskscheduler_delete("reportes_mensuales")  # para borrarla
\end{rnoejecutar}

\begin{advertencia}
Una tarea programada se ejecuta \emph{sin ti}: nadie verá los errores en
pantalla. Por eso es indispensable el registro (\archivo{registro\_reportes.csv}
y \archivo{cron.log}) y que el script use rutas que funcionen desde la
carpeta indicada. Además, la computadora debe estar encendida a esa hora; en
empresas, estas tareas suelen vivir en un servidor.
\end{advertencia}

\subsection{Enviar los reportes por correo}

El paquete \paquete{blastula} permite redactar y enviar correos desde R, con
el reporte como adjunto. Las credenciales (usuario y contraseña del correo)
\textbf{nunca} se escriben en el script: se guardan una sola vez en un
archivo cifrado o en el llavero del sistema.

\begin{rnoejecutar}
install.packages("blastula")
library(blastula)

# UNA sola vez, de forma interactiva: guarda las credenciales cifradas
# en el llavero del sistema (te pide la contraseña, no queda en el código)
create_smtp_creds_key(id = "correo_bi", user = "reportes@techretail.mx",
                      host = "smtp.office365.com", port = 587,
                      use_ssl = TRUE)

# En generar_reportes.R, después de generar cada archivo:
correo <- compose_email(
  body = md(paste0(
    "Hola, adjunto el reporte ejecutivo de **", nombre, "** del mes ",
    mes_reporte, ".\n\nÁrea de Business Intelligence"))
) %>%
  add_attachment(file = file.path(carpeta_salida, archivo))

smtp_send(correo,
          to = "gerente.centro@techretail.mx",
          from = "reportes@techretail.mx",
          subject = paste("Reporte de ventas", mes_reporte, "-", nombre),
          credentials = creds_key(id = "correo_bi"))
\end{rnoejecutar}

Para enviar a cada gerente su reporte, agrega a \archivo{tiendas.csv} (o a
una tabla aparte) una columna con el correo del gerente y úsala dentro del
bucle. Otra alternativa segura para las credenciales son las variables de
entorno: una línea \codigo{CORREO\_PASS=...} en el archivo
\archivo{.Renviron} de tu usuario (que no se comparte ni se sube a Git) y
\codigo{Sys.getenv("CORREO\_PASS")} en el script.

\begin{consejo}
Antes de automatizar el envío, deja que el sistema genere los reportes
durante uno o dos meses mientras \emph{tú} los revisas y los envías a mano.
Cuando confíes en ellos, activa el envío automático. Y agrega una regla de
seguridad: si el registro tiene algún \codigo{ERROR}, que el script te envíe
un correo a ti en vez de a los gerentes.
\end{consejo}

% ----------------------------------------------------------------------------
\section{Reportes en Excel con writexl}
\label{sec:m8-excel}

Muchos usuarios de negocio no quieren un PDF: quieren ``el Excel'' para
filtrar, hacer sus propias tablas dinámicas o pegarlo en su presupuesto. En
vez de pelear contra eso, dales un Excel bien hecho, generado
automáticamente. El paquete \paquete{writexl} escribe archivos
\archivo{.xlsx} sin necesitar Excel ni Java, y si le das una \emph{lista}
de data frames, crea una hoja por cada elemento.

\begin{ejemplo}{Libro de Excel de ventas 2023}
Crear \archivo{resultados/reporte\_ventas.xlsx} con cuatro hojas: un
resumen de KPIs, las ventas mensuales con su variación, las ventas por
tienda y el detalle de todas las transacciones con el nombre de la tienda y
del producto. Después, releerlo para comprobar que quedó bien.
\end{ejemplo}

\begin{rcode}
library(writexl)   # escribir .xlsx
library(readxl)    # leer .xlsx (para verificar)

# Hoja "Detalle": transacciones con nombres legibles
detalle <- transacciones %>%
  left_join(tiendas %>% select(tienda_id, nombre_tienda),
            by = "tienda_id") %>%
  left_join(productos %>% select(producto_id, nombre_producto, categoria),
            by = "producto_id") %>%
  select(transaccion_id, fecha, mes, nombre_tienda, nombre_producto,
         categoria, cantidad, precio_venta, total_transaccion,
         metodo_pago) %>%
  arrange(fecha, transaccion_id)

# Hoja "Resumen": un KPI por fila
resumen <- tibble(
  indicador = c("Ventas totales", "Transacciones", "Unidades vendidas",
                "Ticket promedio"),
  valor = c(round(sum(detalle$total_transaccion), 2), nrow(detalle),
            sum(detalle$cantidad),
            round(mean(detalle$total_transaccion), 2))
)

# Hoja "Mensual": ventas por mes y variación contra el mes anterior
mensual <- detalle %>%
  group_by(mes) %>%
  summarise(ventas = round(sum(total_transaccion), 2),
            transacciones = n(), .groups = "drop") %>%
  mutate(variacion_pct = round((ventas / lag(ventas) - 1) * 100, 1))

# Hoja "Por tienda"
por_tienda <- detalle %>%
  group_by(nombre_tienda) %>%
  summarise(ventas = round(sum(total_transaccion), 2),
            transacciones = n(),
            ticket_promedio = round(mean(total_transaccion), 2),
            .groups = "drop") %>%
  arrange(desc(ventas))

resumen
\end{rcode}
\begin{rsalida}
# A tibble: 4 x 2
  indicador           valor
  <chr>               <dbl>
1 Ventas totales    959993.
2 Transacciones       1000 
3 Unidades vendidas   2969 
4 Ticket promedio      960.
\end{rsalida}

Ahora escribimos las cuatro hojas de una vez. Los nombres de la lista son
los nombres de las pestañas (pueden llevar espacios si los pones entre
acentos graves):

\begin{rcode}
dir.create("resultados", showWarnings = FALSE)

write_xlsx(
  list(Resumen      = resumen,
       Mensual      = mensual,
       `Por tienda` = por_tienda,
       Detalle      = detalle),
  path = "resultados/reporte_ventas.xlsx"
)

# Verificación: ¿qué hojas tiene y cuántas filas hay en cada una?
hojas <- excel_sheets("resultados/reporte_ventas.xlsx")
hojas
sapply(hojas, function(h) {
  nrow(read_excel("resultados/reporte_ventas.xlsx", sheet = h))
})
\end{rcode}
\begin{rsalida}
[1] "Resumen"    "Mensual"    "Por tienda" "Detalle"   
   Resumen    Mensual Por tienda    Detalle 
         4         12         10       1000 
\end{rsalida}

Las cuatro hojas están ahí y el número de filas cuadra: 4 KPIs, 12 meses, 10
tiendas y 1,000 transacciones. Releer lo que escribiste es una verificación
barata que te ahorra enviar un archivo vacío. Comprobemos también que las
fechas y números no se convirtieron en texto:

\begin{rcode}
read_excel("resultados/reporte_ventas.xlsx", sheet = "Mensual") %>%
  head(4)
\end{rcode}
\begin{rsalida}
# A tibble: 4 x 4
  mes     ventas transacciones variacion_pct
  <chr>    <dbl>         <dbl>         <dbl>
1 2023-01 82560.            98          NA  
2 2023-02 65899.            79         -20.2
3 2023-03 90931.            91          38  
4 2023-04 57038.            71         -37.3
\end{rsalida}

Los números llegan como \codigo{<dbl>}, así que en Excel se podrán sumar y
graficar. La variación de enero es \codigo{NA} (no hay mes anterior dentro
de los datos) y en Excel aparecerá como celda vacía.

El archivo \archivo{codigo/reportes/reporte\_excel.R} contiene la versión
completa de este ejemplo, con una hoja adicional de los 15 productos más
vendidos. Para ejecutarlo, copia el archivo a tu carpeta
\archivo{reportes/} y usa
\codigo{source("reportes/reporte\_excel.R")}.

\subsection{Formato avanzado con openxlsx}

\paquete{writexl} es rápido y sencillo, pero no da formato (colores, ancho
de columnas, formato de moneda). Si lo necesitas, usa \paquete{openxlsx}:

\begin{rnoejecutar}
install.packages("openxlsx")
library(openxlsx)
libro <- createWorkbook()
addWorksheet(libro, "Por tienda")
writeData(libro, "Por tienda", por_tienda,
          headerStyle = createStyle(textDecoration = "bold",
                                    fgFill = "#173F69",
                                    fontColour = "white"))
addStyle(libro, "Por tienda", createStyle(numFmt = "$#,##0.00"),
         rows = 2:11, cols = c(2, 4), gridExpand = TRUE)
setColWidths(libro, "Por tienda", cols = 1:4, widths = "auto")
freezePane(libro, "Por tienda", firstRow = TRUE)   # fija el encabezado
saveWorkbook(libro, "resultados/reporte_formato.xlsx", overwrite = TRUE)
\end{rnoejecutar}

\begin{hazlotu}
Agrega al libro de Excel una hoja \codigo{"Por categoría"} con las ventas y
unidades por \codigo{categoria} de producto, ordenada de mayor a menor, y
comprueba con \codigo{excel\_sheets()} que ahora hay cinco hojas.
\end{hazlotu}

% ----------------------------------------------------------------------------
\section{Quarto: el sucesor de R Markdown}
\label{sec:m8-quarto}

\termino{Quarto} es la nueva generación de R Markdown, creada por la misma
empresa (Posit, antes RStudio). Funciona con R, pero también con Python y
Julia, y viene incluido en las versiones recientes de RStudio. La buena
noticia: \textbf{casi todo lo que aprendiste en este módulo sirve igual}. El
Markdown es el mismo, knitr sigue ejecutando el código R y pandoc sigue
convirtiendo el documento. Cambian tres cosas:

\begin{enumerate}
  \item La extensión: \archivo{.qmd} en lugar de \archivo{.Rmd}.
  \item En el YAML, \codigo{output: html\_document} pasa a ser
        \codigo{format: html} (y \codigo{pdf}, \codigo{docx}).
  \item Las opciones de los bloques se escriben preferentemente
        \emph{dentro} del bloque, en líneas que empiezan con \codigo{\#|}, y
        con guiones en lugar de puntos (\codigo{fig-cap} en vez de
        \codigo{fig.cap}).
\end{enumerate}

\begin{textocodigo}[title=reportes/reporte-quarto.qmd]
---
title: "Reporte de ventas"
format:
  html:
    toc: true
  docx: default
params:
  tienda: 1
execute:
  echo: false
  warning: false
---

```{r}
#| label: cargar-datos
library(dplyr)
library(readr)
ventas <- read_csv("../datasets/transacciones.csv",
                   show_col_types = FALSE)
ventas_tienda <- filter(ventas, tienda_id == params$tienda)
```

La tienda `r params$tienda` vendió
$`r round(sum(ventas_tienda$total_transaccion))` en 2023.

```{r}
#| label: fig-mensual
#| fig-cap: "Ventas mensuales"
#| fig-width: 7
ventas_tienda %>%
  count(mes, wt = total_transaccion) %>%
  with(barplot(n, names.arg = substr(mes, 6, 7)))
```
\end{textocodigo}

Las opciones globales ya no van en un bloque \codigo{setup}: se ponen en el
YAML, bajo \codigo{execute}. Para generarlo usas el botón \menu{Render} (el
equivalente de Knit) o la terminal:

\begin{textocodigo}
quarto render reportes/reporte-quarto.qmd              # todos los formatos
quarto render reportes/reporte-quarto.qmd --to docx   # solo Word
quarto render reportes/reporte-quarto.qmd -P tienda:3 # con parámetro
\end{textocodigo}

Desde R, el paquete \paquete{quarto} ofrece \codigo{quarto\_render()},
equivalente a \codigo{rmarkdown::render()}:

\begin{rnoejecutar}
install.packages("quarto")
quarto::quarto_render("reportes/reporte-quarto.qmd",
                      output_format = "html",
                      execute_params = list(tienda = 3))
\end{rnoejecutar}

\begin{consejo}
¿R Markdown o Quarto? Si empiezas un proyecto nuevo, Quarto es la apuesta a
futuro, sobre todo si tu equipo también usa Python o quieres hacer
presentaciones y sitios web. Pero R Markdown no va a desaparecer: sigue
mantenido, hay miles de reportes en producción escritos con él y Quarto
puede generar la mayoría de los \archivo{.Rmd} sin cambios. Lo que aprendas
con uno lo aplicas al otro.
\end{consejo}

% ----------------------------------------------------------------------------
\section{Buenas prácticas para reportes automáticos}
\label{sec:m8-buenas-practicas}

Un reporte automático que falla en silencio es peor que un reporte manual:
genera confianza en números equivocados. Estas prácticas lo evitan.

\subsection{Reproducibilidad}

\begin{itemize}
  \item \textbf{Rutas relativas}, nunca \codigo{setwd("C:/Users/andrea/...")}
        dentro del reporte: en otra computadora esa ruta no existe. Usa
        rutas relativas a la carpeta del proyecto, \codigo{knit\_root\_dir}
        o \paquete{here}.
  \item \textbf{El \archivo{.Rmd} debe ser autocontenido}: cada objeto que
        usa lo crea él mismo. Knit se ejecuta en una sesión limpia y
        \emph{no ve} los objetos de tu consola (sección~\ref{sec:m8-errores}).
  \item \textbf{\codigo{set.seed()}} antes de cualquier cálculo aleatorio
        (muestras, simulaciones, k-means del Módulo~\ref{cap:ml}), para que
        dos ejecuciones den el mismo reporte.
  \item \textbf{Versiones de paquetes}: una actualización de un paquete
        puede cambiar un resultado o romper el código. El paquete
        \paquete{renv} guarda las versiones exactas que usa tu proyecto en un
        archivo \archivo{renv.lock}, y en otra computadora
        \codigo{renv::restore()} instala esas mismas versiones.
\end{itemize}

\begin{rnoejecutar}
install.packages("renv")
renv::init()       # crea una biblioteca de paquetes propia del proyecto
renv::snapshot()   # guarda las versiones actuales en renv.lock
renv::restore()    # (en otra PC) instala exactamente esas versiones
\end{rnoejecutar}

\subsection{Organización}

\begin{itemize}
  \item \textbf{Separa datos, código y salidas}: \archivo{datasets/} (entrada,
        no se modifica a mano), \archivo{reportes/} (plantillas y scripts) y
        \archivo{resultados/} (todo lo que se puede borrar y regenerar).
  \item \textbf{Nombres de archivo} sin espacios ni acentos, con fechas en
        formato \archivo{AAAA-MM-DD} (se ordenan solos) y con lo que
        identifica al reporte: \archivo{2023-12\_T03\_sucursal-sur\_2024-01-05.pdf}.
  \item \textbf{Nombra tus chunks} (\codigo{cargar-datos},
        \codigo{grafica-tendencia}): cuando algo falla, el mensaje de error
        dice el nombre del bloque y encuentras el problema en segundos.
  \item \textbf{Pon la lógica pesada en funciones o scripts} y deja en el
        \archivo{.Rmd} la presentación. Si una función se usa en varios
        reportes, guárdala en un \archivo{funciones.R} y cárgala con
        \codigo{source()}.
\end{itemize}

\subsection{Calidad: revisa antes de enviar}

\begin{itemize}
  \item \textbf{Valida los datos de entrada} dentro del reporte y detente con
        un error claro (\codigo{stop()}) si algo no cuadra: mes sin datos,
        tienda inexistente, columnas faltantes.
  \item \textbf{Cuadra los totales}: la suma de las tiendas debe dar el total
        de la empresa; compara un par de cifras con el sistema fuente.
  \item \textbf{Lee el reporte generado}, al menos al principio y cada vez
        que cambies la plantilla: ¿las frases tienen sentido para una tienda
        que cayó y para una que creció? ¿se ve bien en PDF y en Word?
  \item \textbf{Advierte las limitaciones}: si hay pocos datos, dilo en el
        propio reporte; si detectas un problema de calidad (como los
        productos con costo mayor al precio de catálogo que viste en módulos
        anteriores), menciónalo.
\end{itemize}

% ----------------------------------------------------------------------------
\section{Caso práctico integral: reporte ejecutivo mensual automático}
\label{sec:m8-caso}

\begin{negocio}{El encargo}
La directora comercial de TechRetail México pide: ``Quiero que cada gerente
reciba el primer día hábil del mes un reporte de su tienda con los números
del mes anterior, cómo va contra el mes previo, cómo se compara con las
demás tiendas, qué productos vendió más y una conclusión clara de qué hacer.
En PDF, para imprimirlo en la junta; y quiero poder generar el de cualquier
tienda y mes cuando lo necesite''.
\end{negocio}

La solución tiene tres piezas que ya conoces: la plantilla parametrizada
\archivo{02-reporte-ejecutivo.Rmd}, el script \archivo{generar\_reportes.R}
y la tarea programada. Recorramos la plantilla por partes.

\subsection{Parte 1: encabezado y configuración}

\begin{textocodigo}[title=reportes/02-reporte-ejecutivo.Rmd (encabezado)]
---
title: "Reporte ejecutivo de ventas"
subtitle: "Tienda `r params$tienda` | Periodo `r params$mes`"
author: "Área de Business Intelligence - TechRetail México"
date: "Generado el `r format(Sys.Date(), '%d/%m/%Y')`"
lang: es-MX
geometry: margin=2cm
header-includes:
  - \AtBeginDocument{\renewcommand{\tablename}{Tabla}}
params:
  tienda:
    label: "Tienda (1 a 10)"
    value: 1
    input: select
    choices: [1, 2, 3, 4, 5, 6, 7, 8, 9, 10]
  mes:
    label: "Mes del reporte (AAAA-MM)"
    value: "2023-12"
    input: text
  ruta_datos:
    label: "Carpeta de datos"
    value: "datasets"
    input: text
output:
  html_document:
    toc: true
    toc_float: true
    theme: flatly
  pdf_document:
    fig_caption: true
    extra_dependencies: ["float"]
  word_document:
    fig_caption: true
---
\end{textocodigo}

El subtítulo usa los parámetros, así que cada reporte dice de qué tienda y
periodo es. \codigo{lang: es-MX} pone en español las palabras automáticas
(``Figura'') y \codigo{header-includes} agrega una instrucción de LaTeX para
que las tablas del PDF digan ``Tabla'' y no ``Cuadro''. El paquete
\codigo{float} de \codigo{extra\_dependencies} permite fijar cada figura
donde aparece en el texto (opción \codigo{fig.pos = "H"}).

\begin{textocodigo}[title=Bloque configuracion]
```{r configuracion, include=FALSE}
# ---- Opciones globales: un reporte ejecutivo NO muestra código ----
es_word <- knitr::pandoc_to("docx")   # TRUE si la salida es Word
knitr::opts_chunk$set(
  echo = FALSE, message = FALSE, warning = FALSE,
  fig.width = 7, fig.height = 3.4, out.width = "90%",
  fig.pos = "H",                      # PDF: la figura va donde está
  fig.align = if (es_word) "default" else "center"
)
library(dplyr)
library(readr)
library(ggplot2)
library(scales)
library(knitr)

# Colores corporativos
color_principal <- "#2a78d6"
color_resalte   <- "#eb6834"
color_gris      <- "#b5b3ad"
tema_reporte <- theme_minimal(base_size = 11) +
  theme(panel.grid.minor = element_blank(),
        plot.title = element_text(face = "bold"))

# Nombres de los meses en español (no dependen del idioma de la PC)
meses_es <- c("enero", "febrero", "marzo", "abril", "mayo", "junio",
              "julio", "agosto", "septiembre", "octubre", "noviembre",
              "diciembre")
nombre_mes <- function(m) {
  paste(meses_es[as.integer(substr(m, 6, 7))], "de", substr(m, 1, 4))
}
# Formatos de negocio ("s/d" = sin dato cuando el valor es NA)
dinero  <- function(x) ifelse(is.na(x), "s/d", dollar(x, accuracy = 1))
dinero2 <- function(x) ifelse(is.na(x), "s/d", dollar(x, accuracy = 0.01))
pct     <- function(x) ifelse(is.na(x), "s/d", percent(x, accuracy = 0.1))
```
\end{textocodigo}

\subsection{Parte 2: datos, validación y KPIs}

El bloque \codigo{datos} busca la carpeta de datos (como viste en la
sección~\ref{sec:m8-directorio}), lee y une transacciones y productos, valida
los parámetros con \codigo{stop()} (tienda inexistente, mes sin datos),
calcula el mes anterior con el truco del día 1 menos un día y filtra la
tienda en el mes actual y el anterior (\codigo{ventas\_mes},
\codigo{ventas\_previo}). Es el mismo código que probaste en la consola en
la sección~\ref{sec:m8-parametros}.

El bloque \codigo{kpis} calcula los indicadores con una función que se
aplica igual al mes actual y al anterior (así no hay dos cálculos que puedan
desincronizarse), la posición de la tienda en el ranking del mes y la frase
de variación. Antes de ponerlo en la plantilla, probemos la función en la
consola con la tienda 1 en diciembre, el mismo caso que verás en la figura
del reporte:

\begin{rcode}
calcular_kpis <- function(datos) {
  tibble(
    ventas        = sum(datos$total_transaccion),
    transacciones = nrow(datos),
    unidades      = sum(datos$cantidad),
    ticket        = if (nrow(datos) > 0) ventas / transacciones
                    else NA_real_
  )
}

tienda_1 <- transacciones %>% filter(tienda_id == 1)
kpi     <- calcular_kpis(tienda_1 %>% filter(mes == "2023-12"))
kpi_ant <- calcular_kpis(tienda_1 %>% filter(mes == "2023-11"))
bind_rows(diciembre = kpi, noviembre = kpi_ant, .id = "periodo")

# Ranking de las tiendas en diciembre
ranking <- ventas_dic %>% mutate(posicion = row_number())
ranking %>% filter(nombre_tienda == "Sucursal Centro") %>%
  select(nombre_tienda, ventas, posicion)
\end{rcode}
\begin{rsalida}
# A tibble: 2 x 5
  periodo   ventas transacciones unidades ticket
  <chr>      <dbl>         <int>    <dbl>  <dbl>
1 diciembre  5020.             5       16  1004.
2 noviembre 10593.             9       24  1177.
# A tibble: 1 x 3
  nombre_tienda   ventas posicion
  <chr>            <dbl>    <int>
1 Sucursal Centro  5020.        8
\end{rsalida}

Fíjate en el detalle de \codigo{calcular\_kpis()}: dentro de
\codigo{tibble()} la columna \codigo{ticket} usa las columnas
\codigo{ventas} y \codigo{transacciones} que se acaban de crear, y si no hay
ventas devuelve \codigo{NA} en lugar de dividir entre cero. La Sucursal
Centro vendió \$5,020 en 5 transacciones y quedó en el lugar 8 de 10. Estos
son los números que deben aparecer en el reporte: \emph{así se revisa un
reporte automático}, comparando un caso contra un cálculo independiente.

\begin{textocodigo}[title=Bloque kpis (extracto: la frase de variación)]
texto_variacion <- if (!hay_previo) {
  "no se puede comparar porque no hay ventas registradas el mes anterior"
} else if (var_ventas >= 0) {
  paste0("crecieron **", pct(var_ventas), "** respecto a ",
         nombre_mes(mes_previo), " (", dinero(kpi_ant$ventas), ")")
} else {
  paste0("cayeron **", pct(abs(var_ventas)), "** respecto a ",
         nombre_mes(mes_previo), " (", dinero(kpi_ant$ventas), ")")
}
\end{textocodigo}

\subsection{Parte 3: el texto del resumen y la tabla de KPIs}

Con todo calculado, el resumen es un párrafo normal de Markdown salpicado de
código en línea:

\begin{textocodigo}
# `r info_tienda$nombre_tienda` (`r info_tienda$ciudad`)

## Resumen del mes

En **`r nombre_mes(mes_actual)`**, la **`r info_tienda$nombre_tienda`**
vendió **`r dinero(kpi$ventas)`** en `r kpi$transacciones` transacciones,
con un ticket promedio de **`r dinero2(kpi$ticket)`**. Las ventas
`r texto_variacion`. La tienda ocupó el lugar **`r posicion` de
`r n_tiendas`** en ventas del mes y aportó el **`r pct(participa)`** de las
ventas de la empresa.
\end{textocodigo}

Debajo, el bloque \codigo{tabla-kpis} arma un tibble con cuatro filas
(ventas, transacciones, unidades y ticket), las columnas ``Mes actual'',
``Mes anterior'' y ``Variación'' ya formateadas con \codigo{dinero()} y
\codigo{pct()}, y lo muestra con \codigo{kable()} y el título
\codigo{paste("Indicadores clave de", nombre\_mes(mes\_actual))}.

\subsection{Parte 4: las gráficas}

El reporte tiene tres gráficas que siguen las reglas del
Módulo~\ref{cap:visualizacion}: la tendencia del año con el mes del reporte
resaltado en naranja y el resto en gris, el ranking de tiendas con la tienda
del reporte resaltada, y las ventas por categoría en un solo color. Los
títulos de las figuras (\codigo{fig.cap}) son variables de R, de modo que
también dicen el mes correcto:

\begin{textocodigo}[title=Bloque grafica-tendencia]
```{r grafica-tendencia, fig.cap = titulo_tendencia}
mensual <- ventas_tienda %>%
  filter(mes <= mes_actual) %>%
  group_by(mes) %>%
  summarise(ventas = sum(total_transaccion), .groups = "drop") %>%
  mutate(es_actual = mes == mes_actual,
         # "2023-12" -> "dic" (y respeta el orden cronológico)
         etiqueta = substr(meses_es[as.integer(substr(mes, 6, 7))], 1, 3),
         etiqueta = factor(etiqueta, levels = etiqueta))

ggplot(mensual, aes(x = etiqueta, y = ventas, fill = es_actual)) +
  geom_col(width = 0.7) +
  scale_fill_manual(values = c(`FALSE` = color_gris,
                               `TRUE` = color_resalte),
                    guide = "none") +
  scale_y_continuous(labels = label_dollar(scale = 1e-3, suffix = " mil")) +
  labs(x = NULL, y = "Ventas") +
  tema_reporte
```
\end{textocodigo}

El filtro \codigo{mes <= mes\_actual} hace que el reporte de junio muestre
solo enero a junio: un reporte de un mes pasado no debe ``ver el futuro''.
Las otras dos gráficas siguen el mismo patrón (consúltalas en el archivo
completo).

\subsection{Parte 5: conclusiones automáticas}

El último bloque aplica las reglas de negocio de la
sección~\ref{sec:m8-texto} y escribe una lista de Markdown con
\codigo{results = "asis"}:

\begin{textocodigo}[title=Bloque conclusiones (extracto)]
```{r conclusiones, results = "asis"}
conclusiones <- character(0)

# 1. Crecimiento contra el mes anterior
if (!hay_previo) {
  conclusiones <- c(conclusiones,
    "No hay ventas del mes anterior para comparar el crecimiento.")
} else if (var_ventas > 0.10) {
  conclusiones <- c(conclusiones, paste0(
    "**Crecimiento sólido:** las ventas subieron ", pct(var_ventas),
    ". Identifica qué funcionó para replicarlo en otras tiendas."))
} else if (var_ventas >= 0) {
  # ... "Estabilidad" ...
} else if (var_ventas > -0.10) {
  # ... "Ligera caída" ...
} else {
  conclusiones <- c(conclusiones, paste0(
    "**Alerta:** las ventas cayeron ", pct(abs(var_ventas)),
    ". Revisa inventario, personal y promociones con el gerente."))
}
# 2. Posición (top 3 / 3 últimas)   3. Ticket contra la empresa
# 4. Pocos datos: advertir que las variaciones pueden ser azar
if (kpi$transacciones < 10) {
  conclusiones <- c(conclusiones, paste0(
    "Precaución: el mes solo tiene ", kpi$transacciones,
    " transacciones; interpreta las variaciones con cuidado."))
}
# 5. Calidad de datos: productos con costo mayor al precio de catálogo

# Imprimir como lista de Markdown (por eso results = "asis")
cat(paste0("- ", conclusiones), sep = "\n")
```
\end{textocodigo}

La regla 4 es un ejemplo de honestidad analítica: con menos de 10
transacciones en el mes, una caída del 50\,\% puede ser simple azar, y el
reporte lo dice. La regla 5 revisa un problema de calidad de datos que ya
conoces: hay productos cuyo costo es mayor que su precio de catálogo.

\subsection{El resultado}

Generamos el reporte de la Sucursal Centro de diciembre de 2023 en PDF con:

\begin{rnoejecutar}
rmarkdown::render("reportes/02-reporte-ejecutivo.Rmd",
                  output_format = "pdf_document",
                  output_file   = "reporte.pdf",
                  output_dir    = "resultados",
                  params        = list(tienda = 1, mes = "2023-12"),
                  knit_root_dir = getwd())
\end{rnoejecutar}

La figura~\ref{fig:m08-reporte-ejecutivo} muestra las dos primeras páginas
del PDF tal como salieron, sin ningún retoque. Todo lo que ves (el nombre
de la tienda, los montos, la frase ``cayeron 52.6\,\%'', la tabla, la barra
naranja de diciembre) salió del código.

\begin{figure}[H]
\centering
\fbox{\includegraphics[width=0.47\textwidth]{m08-reporte-ejecutivo}}\hfill
\fbox{\includegraphics[width=0.47\textwidth]{m08-reporte-ejecutivo-p2}}
\caption{Páginas 1 y 2 del reporte ejecutivo en PDF generado para la
Sucursal Centro, diciembre de 2023.}
\label{fig:m08-reporte-ejecutivo}
\end{figure}

La tercera página contiene el final de las conclusiones automáticas. Para
esta tienda, el reporte escribió:

\begin{textocodigo}
- Alerta: las ventas cayeron 52.6%. Revisa inventario, personal y
  promociones con el gerente.
- La tienda está entre las 3 con menores ventas (lugar 8 de 10).
- El ticket promedio ($1,003.90) es menor que el de la empresa
  ($1,059.27): hay espacio para venta cruzada.
- Precaución: el mes solo tiene 5 transacciones; interpreta las
  variaciones con cuidado.
\end{textocodigo}

El mismo archivo produce la versión HTML (con índice flotante) y la de Word
(editable), sin cambiar una línea: basta con
\codigo{output\_format = "all"}.

\subsection{Generar los diez reportes}

Con la carpeta del curso como directorio de trabajo, en la terminal
(o en la pestaña \menu{Terminal} de RStudio):

\begin{textocodigo}
Rscript reportes/generar_reportes.R 2023-12 pdf
\end{textocodigo}

\begin{textocodigo}[title=Salida real del script]
Generando 10 reportes del mes 2023-12 en formato pdf
  [ 1/10] Sucursal Centro    OK    2023-12_T01_sucursal-centro_2026-09-23.pdf
  [ 2/10] Sucursal Norte     OK    2023-12_T02_sucursal-norte_2026-09-23.pdf
  [ 3/10] Sucursal Sur       OK    2023-12_T03_sucursal-sur_2026-09-23.pdf
  ...
  [10/10] Sucursal Express   OK    2023-12_T10_sucursal-express_2026-09-23.pdf

Listo: 10 reportes generados, 0 con error. Carpeta: .../resultados/reportes
Registro guardado en: .../resultados/reportes/registro_reportes.csv
\end{textocodigo}

En unos 20 segundos, los diez PDF están en \archivo{resultados/reportes/}.
Si ejecutas el script para un mes sin datos, por ejemplo
\codigo{Rscript reportes/generar\_reportes.R 2024-01}, ningún reporte se
genera, pero el script no se detiene en la primera tienda: recorre las diez,
y el registro explica la causa:

\begin{textocodigo}[title=resultados/reportes/registro\_reportes.csv (extracto)]
fecha_hora,tienda_id,tienda,mes,archivo,estado,detalle,segundos
2026-09-23 22:02:00,1,Sucursal Centro,2023-12,2023-12_T01_sucursal-centro_2026-09-23.pdf,OK,OK,2.2
...
2026-09-23 22:02:28,1,Sucursal Centro,2024-01,2024-01_T01_sucursal-centro_2026-09-23.html,ERROR,ERROR: No hay datos del mes 2024-01,0.4
\end{textocodigo}

\begin{hazlotu}
Genera los reportes de junio de 2023 en Word para todas las tiendas. Abre el
de la tienda que tuvo el mayor crecimiento y comprueba que su conclusión
dice ``Crecimiento sólido''. Luego revisa \archivo{registro\_reportes.csv}:
¿cuántas filas tiene ahora y por qué?
\end{hazlotu}

% ----------------------------------------------------------------------------
\section{Errores comunes}
\label{sec:m8-errores}

Estos son los errores que casi todo el mundo encuentra en sus primeros
reportes. Los mensajes son los que muestra R (con las rutas acortadas).

\begin{advertencia}
\textbf{Archivo de datos no encontrado al hacer Knit.} El código funciona
en la consola, pero al tejer aparece:
\end{advertencia}

\begin{textocodigo}
Quitting from lines 7-9 [cargar] (reporte.Rmd)
Error:
! 'datasets/ventas_retail.csv' does not exist in current working
  directory ('.../curso-bi-rstudio/reportes').
\end{textocodigo}

\textbf{Causa:} el \archivo{.Rmd} se ejecuta en su propia carpeta
(\archivo{reportes/}), no en la del proyecto. \textbf{Solución:} usa
\codigo{knit\_root\_dir} en \codigo{render()} o
\menu{Knit > Knit Directory > Project Directory}, una ruta
\archivo{../datasets/}, el paquete \paquete{here} o la búsqueda robusta con
el parámetro \codigo{ruta\_datos} (sección~\ref{sec:m8-directorio}). El
mensaje te dice cuál es el directorio de trabajo real: léelo.

\begin{advertencia}
\textbf{Objeto no encontrado (pero en la consola sí existe).}
\end{advertencia}

\begin{textocodigo}
Quitting from lines 7-8 [resumen] (reporte.Rmd)
Error:
! objeto 'ventas' no encontrado
\end{textocodigo}

\textbf{Causa:} Knit ejecuta el documento en una \emph{sesión nueva y
limpia}. El objeto \codigo{ventas} existe en tu consola porque lo creaste
ahí, pero el \archivo{.Rmd} nunca lo crea. \textbf{Solución:} el
\archivo{.Rmd} debe leer sus datos y cargar sus paquetes (\codigo{library()})
él mismo. Prueba: reinicia R (\menu{Session > Restart R}) y ejecuta todos
los bloques en orden con \menu{Run > Run All}; si falla ahí, fallará al
tejer. Lo mismo ocurre con \codigo{could not find function "\%>\%"}: falta
un \codigo{library(dplyr)} dentro del documento.

\begin{advertencia}
\textbf{Falta LaTeX para generar PDF.}
\end{advertencia}

\begin{textocodigo}
! sh: 1: pdflatex: not found

Error: LaTeX failed to compile reporte.tex. See
https://yihui.org/tinytex/r/#debugging for debugging tips.
\end{textocodigo}

En Windows el mensaje suele ser \emph{``No LaTeX installation detected
(LaTeX is required to create PDF output)''}. \textbf{Solución:}
\codigo{tinytex::install\_tinytex()} y reinicia RStudio. Si ya tienes
TinyTeX y el error menciona un archivo \archivo{.sty} faltante, ejecuta de
nuevo el render con Internet para que se descargue.

\begin{advertencia}
\textbf{Nombre de chunk duplicado.}
\end{advertencia}

\begin{textocodigo}
Error in parse_block(g[-1], g[1], params.src, markdown_mode) :
  Duplicate chunk label 'grafica', which has been used for the chunk:
plot(1:10)
\end{textocodigo}

\textbf{Causa:} dos bloques tienen el mismo nombre (pasa mucho al copiar y
pegar un bloque). \textbf{Solución:} cada nombre debe ser único; renombra el
segundo (\codigo{grafica-tiendas}).

\begin{advertencia}
\textbf{Error de YAML por sangría.}
\end{advertencia}

\begin{textocodigo}
Error in yaml::yaml.load(..., eval.expr = TRUE) :
  Scanner error: mapping values are not allowed in this context at
  line 5, column 11
\end{textocodigo}

El encabezado que lo produjo tenía un espacio de más antes de
\codigo{theme}:

\begin{textocodigo}
output:
  html_document:
    toc: true
     theme: flatly
\end{textocodigo}

\textbf{Solución:} alinea las opciones del mismo nivel exactamente con la
misma sangría (múltiplos de dos espacios, sin tabuladores). El mensaje
indica la línea y columna \emph{dentro del YAML}. Otros errores de YAML
frecuentes: olvidar los dos puntos, o escribir un título con dos puntos sin
comillas (\codigo{title: Ventas: diciembre}).

\begin{advertencia}
\textbf{Modificar un parámetro.} \codigo{params\$tienda <- 5} dentro del
documento produce \emph{``no se puede cambiar el valor de un vínculo
bloqueado para `params'\,''}. Copia el parámetro a otra variable.
\end{advertencia}

\begin{table}[H]
\centering
\caption{Otros problemas frecuentes y su solución}
\label{tab:m8-otros-errores}
\small
\begin{tabularx}{\textwidth}{X X}
\toprule
\textbf{Síntoma} & \textbf{Solución} \\
\midrule
El reporte sale con mensajes de carga de paquetes y advertencias &
  \codigo{message = FALSE, warning = FALSE} en \codigo{opts\_chunk\$set()} \\
Cada reporte del bucle sobrescribe al anterior & Da un
  \codigo{output\_file} distinto en cada vuelta \\
La tabla sale como texto de consola (con \codigo{\#\#}) & Usa
  \codigo{kable()} en lugar de imprimir el data frame \\
El texto de \codigo{cat()} aparece en un recuadro gris & Agrega
  \codigo{results = "asis"} al bloque \\
Los meses salen en inglés & Usa un vector propio con los nombres en
  español \\
Los números de un reporte ``se parecen'' a los de otra tienda & Usa
  \codigo{envir = new.env()} en \codigo{render()} \\
\bottomrule
\end{tabularx}
\end{table}

% ----------------------------------------------------------------------------
\begin{resumen}
\begin{itemize}
  \item Un \textbf{reporte reproducible} se genera por completo desde los
        datos y el código: sin copiar y pegar, sin errores de transcripción,
        y se actualiza con un clic.
  \item Un \archivo{.Rmd} tiene \textbf{encabezado YAML}, \textbf{texto en
        Markdown}, \textbf{chunks} de código y \textbf{código en línea}.
        \paquete{knitr} ejecuta el código y \textbf{pandoc} lo convierte a
        HTML, PDF (vía LaTeX) o Word.
  \item Las \textbf{opciones de chunk} controlan qué se ve:
        \codigo{echo}, \codigo{include}, \codigo{message},
        \codigo{fig.cap}, \codigo{results = "asis"}; se fijan para todo el
        documento con \codigo{knitr::opts\_chunk\$set()}.
  \item \codigo{kable()} convierte data frames en tablas de documento;
        \paquete{scales} da formato de negocio a números y porcentajes.
  \item El \textbf{texto dinámico} y las \textbf{conclusiones automáticas}
        (\codigo{if}/\codigo{else}, \codigo{case\_when()}) hacen que la
        narrativa cambie con los datos.
  \item Los \textbf{parámetros} (\codigo{params}) convierten un
        \archivo{.Rmd} en una plantilla; \codigo{render()} con
        \codigo{params}, \codigo{output\_file} y \codigo{knit\_root\_dir}
        genera cada versión.
  \item Un script con un bucle, \codigo{tryCatch()} y un registro genera
        todos los reportes; \textbf{cron} o el \textbf{Programador de tareas}
        lo ejecutan solos, y \paquete{blastula} los envía por correo.
  \item \paquete{writexl} crea libros de Excel con varias hojas a partir de
        una lista de data frames. \textbf{Quarto} (\archivo{.qmd}) es el
        sucesor de R Markdown y usa casi la misma sintaxis.
\end{itemize}
\end{resumen}

\begin{table}[H]
\centering
\caption{Funciones y opciones clave del Módulo 8}
\label{tab:m8-funciones}
\small
\begin{tabularx}{\textwidth}{l X}
\toprule
\textbf{Función u opción} & \textbf{Para qué sirve} \\
\midrule
\codigo{rmarkdown::render()} & Genera el reporte desde código (lo que hace
  el botón Knit) \\
\codigo{params = list(...)} & Valores de los parámetros del reporte \\
\codigo{knit\_root\_dir}, \codigo{output\_dir}, \codigo{output\_file} &
  Carpeta de ejecución, carpeta y nombre de salida \\
\codigo{knitr::opts\_chunk\$set()} & Opciones por defecto de todos los
  chunks \\
\codigo{echo}, \codigo{include}, \codigo{eval}, \codigo{message},
  \codigo{warning} & Qué se muestra y qué se ejecuta \\
\codigo{fig.cap}, \codigo{fig.width}, \codigo{fig.height},
  \codigo{out.width} & Título y tamaño de las figuras \\
\codigo{results = "asis"} & El texto impreso se trata como Markdown \\
\codigo{knitr::kable()} & Tabla con formato de documento \\
\codigo{knitr::knit()} & Teje un documento (o un texto) a Markdown \\
\codigo{knitr::pandoc\_to()} & Formato de salida que se está generando \\
\codigo{dollar()}, \codigo{percent()}, \codigo{comma()} & Formato de
  números en texto y tablas \\
\codigo{tryCatch()} & Atrapa errores para que el bucle continúe \\
\codigo{sprintf()} & Arma nombres de archivo con un molde \\
\codigo{expand\_grid()} & Todas las combinaciones de parámetros \\
\codigo{commandArgs()} & Lee los argumentos de \codigo{Rscript} \\
\codigo{writexl::write\_xlsx()} & Excel con una hoja por data frame \\
\codigo{readxl::excel\_sheets()} & Nombres de las hojas de un Excel \\
\codigo{tinytex::install\_tinytex()} & Instala LaTeX para generar PDF \\
\bottomrule
\end{tabularx}
\end{table}

% ----------------------------------------------------------------------------
\section{Ejercicios}
\label{sec:m8-ejercicios}

Para los ejercicios que piden un archivo \archivo{.Rmd}, guárdalo en la
carpeta \archivo{reportes/} del curso y genéralo con \menu{Knit} o con
\codigo{rmarkdown::render()}. Los demás se resuelven con código R en la
consola, con la carpeta del curso como directorio de trabajo.

\begin{ejercicio}{m8-1}{Tu primer reporte de clientes}
\nivel{Básico}
Crea \archivo{reportes/ej1-clientes.Rmd} con salida HTML que:
(a) tenga título, tu nombre y la fecha automática;
(b) lea \archivo{clientes.csv} en un bloque que no se muestre en el
reporte;
(c) diga en una frase, con código en línea, cuántos clientes hay y cuántos
son premium;
(d) muestre una tabla con \codigo{kable()} del número de clientes por
\codigo{ciudad}, ordenada de mayor a menor.
\end{ejercicio}

\begin{ejercicio}{m8-2}{Tabla ejecutiva de categorías}
\nivel{Básico}
Con \archivo{transacciones.csv} y \archivo{productos.csv}, construye una
tabla con \codigo{kable(format = "pipe")} de las ventas de 2023 por
\codigo{categoria}: ventas en pesos sin centavos, unidades con separador de
miles y porcentaje del total con un decimal, ordenada de mayor a menor y con
el título ``Ventas 2023 por categoría''. ¿Qué categoría vende más y qué
porcentaje representa?
\end{ejercicio}

\begin{ejercicio}{m8-3}{Frase de cumplimiento de meta}
\nivel{Intermedio}
Escribe una función \codigo{frase\_meta(ventas, meta)} que devuelva:
``Se superó la meta en X\,\% (ventas de \$A contra meta de \$B).'' si
\codigo{ventas >= meta}; ``Se alcanzó el X\,\% de la meta\ldots'' si se
alcanzó al menos el 90\,\%; y ``Alerta: solo se alcanzó el X\,\% de la
meta\ldots'' en otro caso. Pruébala con los tres casos y aplícala a las
ventas trimestrales de 2023 (de \archivo{transacciones.csv}) con una meta de
\$220,000 por trimestre.
\end{ejercicio}

\begin{ejercicio}{m8-4}{Opciones de chunk para un reporte ejecutivo}
\nivel{Intermedio}
Crea \archivo{reportes/ej4-opciones.Rmd} con un bloque de configuración
global que oculte código, mensajes y advertencias. Incluye: un bloque
invisible que calcule las ventas por mes de \archivo{ventas\_retail.csv};
una gráfica de líneas con \codigo{fig.cap}, 7 $\times$ 3 pulgadas y
\codigo{out.width = "80\%"}; y un bloque con \codigo{results = "asis"} que
escriba como lista de Markdown los tres meses con más ventas. Genéralo en
HTML y en Word.
\end{ejercicio}

\begin{ejercicio}{m8-5}{Reporte parametrizado por categoría}
\nivel{Intermedio}
Crea \archivo{reportes/ej5-categoria.Rmd} con un parámetro
\codigo{categoria} (valor por defecto \codigo{"Electrónica"}) que muestre,
para esa categoría: ventas totales de 2023 (en una frase), una tabla de sus
5 productos más vendidos y una gráfica de ventas por mes. Después genera,
con \codigo{render()}, los reportes de \codigo{"Ropa"} y \codigo{"Hogar"}
con nombres de archivo distintos en \archivo{resultados/}.
\end{ejercicio}

\begin{ejercicio}{m8-6}{Plan de generación por trimestre}
\nivel{Avanzado}
Construye con \codigo{expand\_grid()} la tabla de trabajos para generar un
reporte por tienda y por trimestre (\codigo{Q1} a \codigo{Q4}). Agrega una
columna \codigo{archivo} con el formato
\archivo{2023-Q1\_T01\_sucursal-centro.pdf} y una columna
\codigo{n\_transacciones} con cuántas transacciones tiene cada combinación
(según \archivo{transacciones.csv}). ¿Cuántos reportes hay que generar?
¿Qué combinación tiene menos transacciones? Marca con
\codigo{generar = FALSE} las que tengan menos de 20.
\end{ejercicio}

\begin{ejercicio}{m8-7}{Un Excel con una hoja por tienda}
\nivel{Avanzado}
Genera \archivo{resultados/ventas\_por\_tienda.xlsx} con una hoja
\codigo{"Resumen"} (ventas y transacciones por tienda) seguida de una hoja
por tienda con el detalle de sus transacciones, cuyo nombre sea el de la
tienda sin la palabra ``Sucursal'' (por ejemplo, \codigo{"Centro"}). Pista:
\codigo{split()} divide un data frame en una lista. Verifica con
\codigo{excel\_sheets()} que hay 11 hojas y que la suma de filas de las
hojas de tiendas es 1,000.
\end{ejercicio}

\begin{ejercicio}{m8-8}{Sistema de alertas con registro}
\nivel{Reto}
Escribe una función \codigo{revisar\_tienda(id, mes)} que calcule la
variación de ventas de una tienda contra el mes anterior y devuelva un
tibble de una fila con la tienda, el mes, la variación y un estado
(``Crecimiento sólido'', ``Estable'', ``Ligera caída'', ``Alerta''), y que
lance un error con \codigo{stop()} si la tienda o el mes no existen, o si el
mes anterior no tiene datos. Recorre todas las tiendas para
\codigo{"2023-06"}, más las combinaciones inválidas \codigo{(11, "2023-06")}
y \codigo{(1, "2023-01")}, atrapando los errores con \codigo{tryCatch()}.
Muestra el registro y escribe con \codigo{cat()} el texto Markdown de
alertas que iría en el reporte.
\end{ejercicio}
